\documentclass[11pt,a4paper]{article}
\usepackage{epsf,amsfonts,amssymb,epsfig,color,fancyhdr,lastpage,mathptmx}
\usepackage{wrapfig}
\PassOptionsToPackage{hyphens}{url}
\usepackage[colorlinks=true,linktocpage=true,linkcolor=blue,citecolor=blue]{hyperref}
%\usepackage[colorlinks=true,linktocpage=true,linkcolor=blue,citecolor=blue]{hyperref}\pagestyle{fancy}
\usepackage[gen]{eurosym}
 \usepackage{enumitem}


\headsep 0.45cm
\topmargin -1cm
\textheight 25.5cm
\oddsidemargin -0.14cm
\textwidth 16.6cm
\linespread{0.97}
\fancyhead{}
\fancyfoot{}
\headwidth 17cm
\usepackage{caption}

\usepackage{longtable}

%\fancyhead[RO,RE]{\small CV Blas}
%\fancyfoot[CO,CE]{\small Page {\thepage} of \pageref{LastPage}}

%\usepackage{cmbright}
%\usepackage[T1]{fontenc}

\usepackage{newtxtext}
%\urlstyle{same}
\def\UrlFont{\tt\small}

%\renewcommand{\rmdefault}{phv} % Arial
%\renewcommand{\sfdefault}{phv} % Arial

%\documentclass[a4paper,10pt,oneside]{article}
%%\usepackage{babel}
%\usepackage[ansinew]{inputenc}


\newcommand{\db}[1]{\textcolor{cyan}{\it{\textbf{db:#1}}}}
%\usepackage{setspace}
%\setlength\itemsep{0.9in}
%\pagestyle{headings}

\newcounter{num}
 \setcounter{num}{1}
\newcommand{\Num}{[\arabic{num}] \addtocounter{num}{1}}

\definecolor{midgreen}{rgb}{0.0,0.675,0.0}

% \newcommand{\mz}[1]{\textcolor{midgreen}{#1}}
% \newcommand{\mzc}[1]{\textcolor{midgreen}{\textbf{[MZ: #1]}}}

\newcommand{\mz}[1]{#1}
\newcommand{\mzc}[1]{}

\usepackage{sectsty}
\sectionfont{\fontsize{12}{15}\selectfont}

\usepackage[gen]{eurosym}

\usepackage{xcolor}
\usepackage{titlesec}

% Define custom colors
\definecolor{sectioncolor}{RGB}{120,0,120}  % Orange
\definecolor{subsectioncolor}{RGB}{80,0,80}  % Purple
% \subsubsection
\definecolor{subsubsectioncolor}{RGB}{100,0,100}  % Purple
\definecolor{paragraphcolor}{RGB}{100,0,100}  % Darker Purple
\definecolor{subparagraphcolor}{RGB}{140,0,140}  % Darker Purple


% Apply colors to sections, subsections, and paragraphs
\titleformat{\section}
{\normalfont\Large\bfseries\color{sectioncolor}}{\thesection}{1em}{}
%left, before after
\titlespacing*{\section}{0pt}{10pt}{7pt}

\titleformat{\subsection}
{\large \normalfont\bfseries\color{subsectioncolor}}{\thesubsection}{1em}{}
\titlespacing*{\subsection}{0pt}{8pt}{5pt}


\titleformat{\subsubsection}
{\normalfont\large\bfseries\color{subsubsectioncolor}}{\thesubsection}{1em}{}

\titleformat{\paragraph}[runin]
{\normalfont\normalsize\bfseries\color{paragraphcolor}}{\theparagraph}{1em}{}
%left, before after
\titlespacing*{\paragraph} {0pt}{5pt}{7pt}

\titleformat{\subparagraph}[runin]
{\normalfont\normalsize\bfseries\color{subparagraphcolor}}{\thesubparagraph}{1em}{}


%CHECK OUT THE FONT! https://www.ipgp.fr/~moguilny/LaTeX/fontawesome5Icons.pdf
\usepackage{fontawesome5}


\begin{document}


\hspace{-.9cm}
\begin{tabular}{l}
{\LARGE{\bf \color{sectioncolor} MIGUEL ZUMALACARREGUI PEREZ}}\\
{}\\[.3cm]
ORCID Number: \href{https://orcid.org/0000-0002-9943-6490}{0000-0002-9943-6490}\\
ADS/NASA: \href{https://ui.adsabs.harvard.edu/search/fq=%7B!type%3Daqp%20v%3D%24fq_database%7D&fq_database=(database%3Aastronomy%20OR%20database%3Aphysics)&q=author%3A%22Zumalacarregui%2C%20Miguel%22&sort=date%20desc%2C%20bibcode%20desc&p_=0}{ui.adsabs.harvard.edu}\\
Google Scholar: \href{https://scholar.google.com/citations?user=Zsvsq0MAAAAJ&hl=en}{scholar.google.com}\\
INSPIRE-HEP: \href{https://inspirehep.net/authors/1113379?ui-citation-summary=true}{inspirehep.net/authors/1113379}
%\href{https://oneresearchcommunity.com/author/orcid-0000-0002-9943-6490}{One score:} 753

%\href{https://scholar.google.com/citations?user=YBqSciQAAAAJ&hl=en#d=gsc_md_hist&t=1655709380773}{Google Scholar}
\end{tabular}
%\newpage
%\vspace{.5cm}


%\vspace{-0.55cm}
%\begin{center}

%%%%%%%%%%%%%%%%%%%%%%%%%%%%%%%%%
%\section*

\vspace{-2cm}

\hspace{8.65cm}
{\bf{Current Position}} %Professional Situation
%%%%%%%%%%%%%%%%%%%%%%%%%%%%%%%%%
\vspace{-0.01cm}

\hspace{8.7cm} Principal Investigator, GLOW (ERC)

% since 01/02/2018\\
\vspace{-0.05cm}

\hspace{8.7cm} Ram\'on y Cajal Fellow
\vspace{-0.05cm}

\hspace{8.7cm} Instituto de F\'isica Te\'orica IFT-UAM/CSIC \vspace{-0.05cm}

\hspace{8.7cm} Madrid, Spain
%\emph{Telephone}: +44 20 7848 2578\\

%%%%%%%%%%%%%%%%%%%%%%%%%%%%%%%%%%
\section*{Personal and Contact Information}
%%%%%%%%%%%%%%%%%%%%%%%%%%%%%%%%%%
% \emph{Date and Place of Birth}: July 27, 1983, Madrid, Spain.\hspace{1cm}\emph{DNI}: 50890256L\\
% \emph{Email}: \href{mailto:mzuma@aei.mpg.de}{mzuma@aei.mpg.de} \hspace{0.3cm} \\
\emph{Webs:} \url{https://miguelzumalacarregui.es/} \\
${}$ ~~~~~~~~~\,\url{https://glow-astro.org/} \\
${}$ ~~~~~~~~~\,\url{https://www.aei.mpg.de/person/104068/337435}
%\emph{Citizenship}: Spanish. \hspace{1cm} \emph{Gender}: Male.\\
  %\emph{Passport Number}: AAD599957.\\
\\
\textit{Language skills:} Spanish (native), English (C2), German (C1)
 
%\emph{Unesco Codes}: 221212, 221214, 221207, 221205.
% \vspace{-0.55cm}

\section*{Summary}

I am a theoretical physicist working at the intersection of cosmology, gravitation, and fundamental physics. I am an expert in gravitational waves, gravitational lensing, dark matter, theories of gravity, dark energy, and interpreting cosmological and gravitational-wave observations. % and am well-versed in data analysis, astrophysics, early-universe cosmology, and high-energy physics.
I strive to push the boundaries between different fields, reaching data-driven conclusions, finding novel applications of existing methods, and combining theoretical and computational approaches.


%%%%%%%%%%%%%%%%%%%%%%%%%%%%%%%%%
\section*{Professional Activity}
%%%%%%%%%%%%%%%%%%%%%%%%%%%%%%%%%
\vspace{-0.3cm}
% \subsection*{\sc Positions Held}\vspace{-0.3cm}
\begin{longtable}{c p{15.5 cm}}
2026- &
\textbf{Principal Investigator} of the ERC Consolidator Grant \textbf{GLOW} (\href{https://glow-astro.org/}{glow-astro.org}) \\ & \& \textbf{Ram\'on y Cajal Fellow} at the \textbf{Instituto de F\'isica Te\'orica IFT-UAM/CSIC}, Madrid (Spain). \\[3pt]
2020- & 
\textbf{Group Leader} at the \textbf{Max Planck Institute for Gravitational Physics} \\ & (Albert Einstein Institute), {Astrophysical \& Cosmological Relativity division}, Potsdam (Germany). \\[3pt]
% \textbf{Junior Group Leader} in the \textit{Astrophysical and Cosmological Relativity division}, \\ & \textbf{Max Planck Institute for Gravitational Physics} (Albert Einstein Institute) Potsdam (Germany). \\[3pt]
2017-20 & \textbf{Marie Curie Global Fellow} at the \textbf{Berkeley Center for Cosmological Physics}, Berkeley (US) \\ & with return phase at the \textbf{Institut de Physique Th\'eorique} CEA, Saclay (France). 
% \\ & Hosts: Uro\v{s} Seljak (Berkeley) and Filippo Vernizzi (Saclay). %[Awarded 2016].
\\[3pt]
2015-17 & \textbf{Nordita Fellow} at the \textbf{Nordic Institute for Theoretical Physics}, Stockholm (Sweden).
\\[3pt]
2013-15 & Postdoctoral  Researcher at \textbf{Institut f\"ur Theoretische Physik - Uni. Heidelberg} (Germany). 
\\[3pt]
2013 & Postdoctoral Researcher at \textbf{Instituto de F\'isica Te\'orica IFT-UAM/CSIC}
(Madrid, Spain).
\\[3pt]
2008-12 & Graduate student (FPI) at the \textbf{Institut de Ciencies del Cosmos - University of Barcelona} (ICC-UB) and the \textbf{IFT-UAM/CSIC} (Spain) 
\end{longtable}

%%%%%%%%%%%%%%%%%%%%%%%%%%%%%%%%%
\section*{Academic Degrees}
%%%%%%%%%%%%%%%%%%%%%%%%%%%%%%%%%
\vspace{-0.3cm}
\begin{longtable}{c p{15.5 cm}}
2012 & \textbf{Doctor in Physics}, \emph{Cum Laude} \& International Mention, Universidad Aut\'onoma de Madrid
\\ &
 Dissertation: \emph{``Probing the foundations of the standard cosmological model''} 
% \\ & Advisors: Juan Garcia-Bell\'ido, Pilar Ruiz-Lapuente and Tomi S. Koivisto.  
\\[3pt]
2009 & \textbf{Master in Theoretical Physics}, Universidad Aut\'onoma de Madrid (120 ECTS) 
\\ &
 Thesis: \emph{``Disformal Quintessence: observational constraints and linear perturbations''} 
\\[3pt]
2007 & \textbf{Physics Degree}, Universidad Aut\'onoma de Madrid (300 ECTS)
\end{longtable}

% \section*{International Projection}
%%%%%%%%%%%%%%%%%%%%%%%%%%%%%%%%%


%%%%%%%%%%%%%%%%%%%%%%%%%%%%%%%%%
\section*{Career disruptions \& major life events}
%%%%%%%%%%%%%%%%%%%%%%%%%%%%%%%%%
% \vspace{-0.3cm}
% [You may include a short factual explanation of career breaks or diverse career paths such as secondments, volunteering, part-time work, time spent in different sectors or the effects of major life events such as long term illness as well as the effects of pandemic restrictions on research productivity.]
\begin{itemize}\setlength{\itemsep}{1pt}
\item Two children born in 2019 and 2022. I have officially taken \textbf{19 weeks of parental leave}. % (shifting my PhD year to 2013).
\item Compounding disruptions during the \textbf{COVID-19 pandemic}, including lack of childcare or family support for 2 working parents and international move from California to Germany in Aug. 2020. 
% \item My spouse had a serious health crisis during her second pregnancy in 2021, during which I had to take care of her and my son 6 months. 
\item In 2021 my son showed strong signs of autism and was diagnosed in 2022. I accompany him to regular therapies and navigate support and integration in the German education system.
% Since then I share the duty of accompanying my son to regular therapies. %Three changes in daycare.
\item Family reasons force me to \textbf{decline numerous invitations} for seminars, conferences and review panels. % (including 3 plenary talks at conferences in 2023 alone) and often disrupt my work.
\end{itemize}


% The 2019-22 disruptions coincided with \textbf{establishing a new group and embarking on a new direction} (GW lensing) and prevented me from applying for substantial funding schemes, except for supporting candidates for external fellowships.
% Since 2022, I have shared the duty of accompanying my son to therapy several days a week.
% These transient setbacks have affected my production, but they have also made me a more resourceful scientist and a kinder person.
% % The daycare closed unexpectedly the entire week before the ERC deadline, also while our babysitter was on vacation (I apologize for typos).
% Family reasons have forced me to \textbf{reject numerous invitations} (including 3 plenary talks at conferences in 2023 alone) and often disrupt my work.
% % (for a glimpse into my life, daycare closed the entire week week before the ERC-CoG deadline, right when the babysitter was on vacation). 
% I have officially taken \textbf{19 weeks of parental leave} (shifting my PhD year to 2013).

%%%%%%%%%%%%%%%%%%%%%%%%%%%%%%%%%
%rtaciones Cientifico-Tecnicas: Aportaciones científico-técnicas: se valorará la relevancia de la contribución de la persona candidata en los artículos publicados en revistas científicas, los libros o capítulos de libros científicos y técnicos, los trabajos presentados en congresos, las patentes concedidas o licenciadas, y, en general, en cualquier otra aportación que permita valorar los diferentes aspectos de la investigación, incluyendo la transferencia de tecnología. Puntuación: de 0 a 10 puntos.

% \section{Other qualifications}

% \subsection*{Languages}

% % \mzc{Brief summary, languages, places worked...}

% Spanish (native), English (C2), German (C1)

\newpage

\section*{Scientific \& technical contributions}
%%%%%%%%%%%%%%%%%%%%%%%%%%%%%%%%%

\subsection*{Publications Summary}


\def\signhot{{\color{gray!30!black}\faHotjar}}
\def\signcited{\faTrophy}
% \def\signsenior{\includegraphics[width=8pt]{images/person-cane-solid.pdf}}%\faUserGraduate} %\faPersonCane
\def\signsenior{\faHatWizard}
\def\signcode{} %{\faCode}
\def\signidea{\faLightbulb[regular]}
\def\signlead{\faFlag[regular]}
\def\signalpha{\faSortAlphaDown}
\def\signfeatured{\faStar[regular]}

\begin{minipage}{0.55\textwidth}

\begin{itemize}%\setlengh\itemsep{2pt}
% {\small \url{http://inspirehep.net/search?p=author\%3AM.Zumalacarregui.1}}
 \item Total: 11,132 citations (average of 146.5 per entry), \\ 76 entries, h-index 42 (\href{http://inspirehep.net/author/profile/M.Zumalacarregui.1}{InSPIRE})
  \item Published $\leq 10$ authors: 5,060 citations (average 103.3/entry), 49 entries, h-index 33 (\href{https://inspirehep.net/literature?sort=mostrecent&size=25&page=1&q=a%20zumalacarregui&author_count=10%20authors%20or%20less&ui-citation-summary=true}{InSPIRE})

 \item {7$\times$ Top 1\% cited in Physics %\signcited 
 \\ 1$\times$ Top 0.1\% \& $\leq 2$ years %\quad \signhot{} 
 \; {(\href{https://www.webofscience.com/wos/woscc/summary/6ce91b13-02c0-414c-ac15-c4f89814dc78-fc3901c0/relevance/1}{Web of Science})}}
 
 % \item Topcite: 6$\times$ 500+, 5$\times$ 250-499, 12$\times$ 100-249, 14$\times$ 50-99 citations.
 % Topcite 37/76 $\approx$ 49\% of entries (50+ citations on \href{http://inspirehep.net/author/profile/M.Zumalacarregui.1}{InSPIRE}), 08 Sep 2026.

 % Counted from InSPIRE (a M.Zumalacarregui.1), 08 Sep 2026, over all 76 entries:
 %   first/single  first author of a non-alphabetical list, or sole author (the
 %                 2018 "No LIGO Macho" reply is a comment and is not counted)
 %   first tier    non-alphabetical, second or third author and not last
 %   alphabetical  author list in alphabetical order, so position carries no meaning
 \item 10$\times$ first/single author, 8$\times$ first tier, 25$\times$ alphabetical \\
 \textbf{Alphabetical ordering is default} in the field. 
\end{itemize}
% {\footnotesize \noindent
% Updated publication list \href{http://inspirehep.net/search?p=author\%3AM.Zumalacarregui.1}{http://inspirehep.net/search?p=author\%3AM.Zumalacarregui.1} - citation count from Inspire-HEP}
\end{minipage}
\begin{minipage}{0.45\textwidth}
\vspace*{-1cm}
\includegraphics[width=0.99\textwidth]{images/metrics_indices_266.png}
\end{minipage}

\vspace{.5cm}
Publications: % list in page \pageref{sec:publication_list},
% Updated lists: 
\href{https://inspirehep.net/search?p=author\%3AM.Zumalacarregui.1}{InSPIRE},
\href{http://adsabs.harvard.edu/cgi-bin/basic_connect?qsearch=miguel+zumalacarregui&version=1}{ADS}, 
\href{https://scholar.google.com/citations?user=Zsvsq0MAAAAJ&hl=en}{Google Scholar},
\href{https://www.webofscience.com/wos/woscc/summary/c4236c65-f24f-4c7f-bba1-97eb6a071a0a-fc37ee8a/relevance/1}{Web of Science}.
% {\small \url{http://miguelzuma.github.io/research.html}.}

\mzc{update plot}


\subsection*{Public Software Tools}
\begin{itemize}
\item \texttt{hi\_class} (Horndeski in the Cosmic Linear Anisotropy Solving System) \quad  \href{https://hiclass-code.net}{hiclass-code.net} \\ Cosmological solver for dark energy \& gravity theories 
\quad $\longrightarrow$ \underline{Used in 99 publications} (\href{https://hiclass-code.net/#pubs}{list})
\item \texttt{GLoW} (Gravitational Lensing of Waves)  \quad \href{https://glow-astro.org/glow-code.html}{glow-astro.org/glow-code} \\ Wave-optics lensing solver for general matter distributions \quad $\longrightarrow$  \underline{Used in 18 publications} (\href{https://github.com/glow-astro/GLoW}{github.com/glow-astro/GLoW})
\end{itemize}
\vspace{-5pt}
I have also contributed to GSHE, Glworia, WOlensing, see \href{https://miguelzumalacarregui.es/software.html}{miguelzumalacarregui.es/software}

\mzc{add Glworia, WOlensing, GSHE, companion papers}

%%%%%%%%%%%%%%%%%%%%%%%%%%%%%%%%%
% Internacionalización: se valorará la participación directa en acciones relacionadas con programas y proyectos internacionales, especialmente relacionadas con programas de movilidad internacional predoctoral y posdoctoral, así como las publicaciones, participación o financiación en proyectos y contratos realizados en colaboración internacional. Puntuación: de 0 a 4 puntos.

\section*{Professional recognition}


\subsection*{Funding}
More than 4 M\euro{} in research projects led directly:  %2.1M
\vspace{-8pt}
\begin{longtable}{l p{14.5 cm}}
2025 & \textbf{ERC Consolidator Grant} - GLOW  \textbf{2,185,007 \euro}. % PE9
\\ & \href{https://www.aei.mpg.de/1359867/erc-consolidator-grant-for-miguel-zumalac-rregui}{\it ``Gravitational Lensing of Waves''} \; (\href{https://glow-astro.org/}{glow-astro.org})
\\[3pt] 
2025 & Ramon y Cajal Fellowship, ``Atraccion de Talento'' \&  RyC MaX \textbf{325,000 \euro}.
% \textbf{100,000} \euro \\[1pt] & CSIC Max bonus \textbf{300,000} \euro 
\\[3pt]
2024 & \textbf{Max Planck Society's Annual Project} (philanthropic donations)
  \textbf{286,940 \euro} (no overhead) %286,940.60 %supporting members  PI
\\[1pt]
& \href{https://www.mpg.de/annualproject2024} {\textit{``With Einstein on crooked paths''}} (50,000 \euro{} \href{https://www.aei.mpg.de/1227766/with-einstein-on-crooked-paths}{above average funding} in the program)
% among the best funded projects in the program's history. %, average 235,000 \euro %(50,000 more)
\\[3pt]
2024 & Invited to interview for ERC Consolidator Grant %- GLOW. % (2,187,171 \euro) %PE9 
% \\ & Score A: ``\textbf{Fully meets the ERC's excellence criterion} and is recommended for funding'' % if sufficient funds are available
\\[3pt]
2024 & ESI workshop ``Lensing and Wave Optics in Strong Gravity'' \textbf{12,000 \euro}\\[3pt]
2020- & \textbf{Group Leader}, Max Planck Institute for Gravitational Physics, Potsdam (Germany) \\ & 5yr funding for PI, 2 postdocs \& 1 PhD from ACR Division, circa \textbf{1,200,000 \euro} (no overhead)\\[3pt]
% 2022-2024 & Host for 2 externally funded postdocs (Humboldt \& Margarita Salas Fellowships) \\[3pt]
2019 &  Invited to interviews for ERC Starting Grant \& ERF %- %DarC-GraM. %(1,497,644 \euro) PE9 
\\[3pt]
2015 & \textbf{Marie Sklowdoska-Curie Global Fellowship} NLO-CO \textbf{264,668 \euro} %\textbf{PI} 
\\[1pt]
& \href{https://cordis.europa.eu/project/id/706434}{\textit{``Large scale structure of the Universe and fundamental physics''}}\\[3pt]
2008 & PhD project: FPI Grant BES-2008-009090 \textbf{54,816 \euro}  %proyect AYA2006-05369 
\\[1pt] & Research stays in Heidelberg \& Oslo \textbf{14,550 \euro}
%\\[1pt] 2012 &  PhD FPI research stay at ITA Oslo (Norway), 5.3 Mo \textbf{8,700 \euro}  % (105 days) %  
% \\[1pt] 2009 & PhD FPI research stay at ITP Heidelberg (Germany), 3.5 Mo \textbf{5,850 \euro}  %(162 days) % 
 \end{longtable}
\vspace{-8pt}\noindent
Host for two awarded Humboldt Fellowships (2022, 2025) and one ``Margarita Salas'' Fellowship (2022).
\\
Senior Collaborator in GWSky ERC Synnergy Grant 12M\euro -- (\href{https://gwsky.org/people/}{link}) 

\subsection*{Leadership in large consortia}
\vspace{-0.3cm}
\begin{longtable}{l p{14.5 cm}}
2025- & \textbf{Senior co-chair - LISA Cosmology} Working Group, with 500+ members.
 \\ & Elected democratically by my peers.  \\[3pt]
2025- & \textbf{Lead - Euclid Gravitational Waves} Science Working Group, with 80+ members
\\ & Appointed by Euclid Consortium Board to establish this new direction. %Deputy
% 2018- & Member of the {LISA Consortium} - \\ 
% % \url{https://www.esa.int/Science_Exploration/Space_Science/LISA}\url{https://www.elisascience.org/} \\
% 2014- & Member of the {Euclid Collaboration} 
% % - \url{http://www.cosmos.esa.int/web/euclid}
\end{longtable}

\newpage

\subsection*{Awards and fellowships}
\begin{itemize}\setlength{\itemsep}{1pt}
    \item \textbf{ERC Consolidator Grant} 2025 -- GLOW, ``Gravitational Lensing of Waves''
    \item R3 certification -- Spanish Research Agency
    \item Ramon y Cajal Fellowship 2024 -- \textbf{Ranked 1 in Astrophysics \& 2 in Physics} 
    \item \textbf{Max Planck Society's Annual Project} 2024 -- ``With Einstein on crooked paths'',
    funded by philanthropic donations
    \item Berkeley Connect Fellow 2018, UC Berkeley (US)
    \item Marie-Sklodowska Curie Global Fellow 2015 - 3yr at UC Berkeley, US \& IPhT Saclay, France
    \item Nordita Fellow 2015 -- 2yr at Nordic Institute for Theoretical Physics, Stockholm (Sweden) %(2015-17) %720,000 SEK -> 74,883 Eur on Oct 2017
    % \item FPI Research stay at ITA Oslo (Norway), 5.3 Months (2012)
    \item Yggdrasil Fellow 2011 -- 2Mo stay in ITA Oslo (Norway) %, 2 Months %34,000 NOK at (0.1227 rate for 15.11.2010) = 4171.8 Eur
    % \item FPI Research stay at ITP Heidelberg (Germany), 3.5 Months (2009)
    \item FPI Fellowship 2008 -- 4yr PhD at ICC 
    \item UAM Master's scholarship 2007 -- 2yr support
    % \item Erasmus program (2005-2006)
    % \item 7 times ``matricula de honor''
\end{itemize}

%%%%%%%%%%%%%%%%%%%%%%%%%%%%%%%%%
% Liderazgo: se valorará la independencia investigadora de la persona candidata mediante su participación como investigador/a principal en programas y proyectos de investigación nacionales, regionales y/o internacionales obtenidos mediante concurrencia competitiva, incluidos los financiados por empresas y otras entidades privadas. Puntuación: de 0 a 3 puntos.
\section*{Scientific evaluation}
%%%%%%%%%%%%%%%%%%%%%%%%%%%%%%%%%


\subsection*{Member of Evaluation Panels}
\begin{itemize}\setlength{\itemsep}{1pt}
\item {NASA Hubble Fellowship Program (NHFP)}, Physics and Cosmology panel (USA). %2022-23, invited the previous year but on parentak leave
\item {FCT} - Portuguese Foundation for Science and Technology Individual Support Call, (Portugal) %2021: 23 applications for junior, associate and principal researchers
\item MPI-Grav - member of committee evaluating $\sim 150$ postdoc applications per year (since 2020) %Artists evaluator
\end{itemize}


\subsection*{External Evaluator for Funding}

\begin{itemize}\setlength{\itemsep}{1pt}
\item {Starting Grants - European Research Council}  
\item {Gravitational Waves Call STFC/UKRI} (United Kingdom)
\item {SNSF} - Swiss National Science Foundation (Switzerland) 
\item {Alexander von Humboldt Foundation} postdoctoral fellowships. 
 \item {InterTalentum - Marie Sklodowska Curie COFUND}
 \item {Netherlands Research Council (NWO)}
 \item Research Grants Council of Hong Kong
 \item {Estonian Research Council}
\end{itemize}

\subsection*{Referee for Publications}
Peer review for 40+ publications in international journals:
\begin{itemize}\setlength{\itemsep}{1pt}
  \item {Physical Review Letters} (14 articles), %last from 09-24 (PTA-H0)
\item {Nature Astronomy} (2 articles), %last from 3-24 (lensing H0)
\item {Physical Review D} (19 articles), %last from 09-26 (type-II LSD)
\item {JCAP} (7 articles), % last from 7-19 (dark sector Horndeski)
\item {The Astrophysical Journal} (5 articles), %Astrophysical Journal, 5 article last from 03-26
\item {Astrophysical Journal Letters} % (1 articles), %APJL 1 article last from 9-25
 \item {Physics Letters B} % (1 evaluated article) %last from ~'12
\item {MNRAS letters} %(1 article, letter) %last from 2-23
\item {Annals of Physics} %(1 evaluated article) %last from ~12-22
 \item {Physics of the Dark Universe} % (2 evaluated articls) % last from 5-18
 \item {European Physical Journal C} %(3 evaluated article) %last from 11-19
\item  {Astroparticle Physics} %(1 evaluated article) %last from 4-18
\end{itemize}
\vspace{-.2cm}
 % {International Journal of Modern Physics D}, %(1 evaluated) %last from 9-17
 % {PASJ}, % Proceedings of the Astronomical Society of Japan (1 evaluated) % last from 7-18 
 % {Entropy}%, %(1 evaluated article) %last from 4-14
%  (1 article each).


\newpage



%%%%%%%%%%%%%%%%%%%%%%%%%%%%%%%%%
%Otros méritos curriculares: se valorará, entre otros aspectos, la dirección de tesis doctorales y trabajos de fin de máster universitarios, la obtención de premios, menciones y distinciones, actividades de divulgación científica y cualquier otra aportación que permita valorar méritos en investigación no incluidos en los anteriores apartados. Puntuación: de 0 a 3 puntos
\section*{Supervision \& coordination}
%%%%%%%%%%%%%%%%%%%%%%%%%%%%%%%%%
Official (co)supervisor to 17 junior scientists (postdocs, PhD, MSc) and informal advisor to 20+ others in ongoing collaborations, research visits and published works.
% Most of my mentees have attained highly competitive fellowships and leadership positions.

\subsection*{Postdoc: official supervision}\setlength{\itemsep}{1pt}
\begin{itemize}\setlength{\itemsep}{1pt}
% \item {Juan Urrutia} (Humboldt Fellow at MPI-Grav, starting 2026)
\item {Xikai Shan} (Postdoc at MPI-Grav, 2026-)
\item {Han-Gil Choi} (Postdoc at MPI-Grav, 2026-)
\item {Srashti Goyal} (Postdoc at MPI-Grav, 2023-2026) $\longrightarrow$ \underline{Independent Fellow} at IUCAA.
\item {Serena Giardino} (Postdoc at MPI-Grav, 2024-2025) %funding from OCR division, MPI-Grav Hannover) 
$\longrightarrow$ \underline{Postdoc at Imperial College} (w/ C. de Rham).
 \item {Hector Villarubia-Rojo} (Margarita Salas Fellow at MPI-Grav 2022-2024) $\longrightarrow$ return phase at UCM.
 \item {Guilherme Brando} (Humboldt Fellow at MPI-Grav, 2022-2024) $\longrightarrow$ \underline{Tenured Researcher at UFdES, Brazil} %Universidade Federal do Espírito Santo
 \item {Giovanni Tambalo} (Postdoc at MPI-Grav, 2020-2023)  $\longrightarrow$ \underline{Postdoc at ETH Zurich} (w/ L. Senatore) \\ $\longrightarrow$ Marie Sklodowska Curie Fellow at IPhT Saclay.
\end{itemize}


\subsection*{Postdocs: informal supervision} \setlength{\itemsep}{1pt}
\begin{itemize}\setlength{\itemsep}{1pt}
% \item {Alberto Mangiagli} (MPI-Grav, 2024-) Diffraction by dark matter in LISA.
\item {Maria Berti} (U. Geneva, 2024-25) Dark energy reconstruction, {1 joint publication}.
% \item {Han-Gil Choi} (IBS, 2024-) visited my group in 2024-25, signatures of dark-matter objects. %2 weeks oct 24, 1 week march/april 25
% \item {Yifan Wang} (MPI-Grav, 2023-) searches for GWs near supermassive black holes.
% \item {Arnab Dhani} (MPI-Grav, 2023-) GW lensing by future ground detectors.
\item {Marius Oancea} (Postdoc at U. Vienna 2022) visited my group in 2021, 2 joint publications
\end{itemize}

\subsection*{PhD thesis supervision}
\begin{itemize}\setlength{\itemsep}{1pt}
 % \item {Francisco Pimienta} PhD in U. Porto 2024-28, co-supervisor.
 
 \item {Nicola Menadeo} PhD at U. Naples, co-supervisor, long-term stay at MPI-Grav 2023-25 \\ $\longrightarrow$ \underline{Postdoc at Humboldt Uni.}

\item {Stefano Savastano} PhD at MPI-Grav, 2020-24, main supervisor (officially w/ A. Buonanno \href{https://edoc.hu-berlin.de/items/19db3995-f584-401e-b2d3-876b5512106b}{thesis link}). %\textbf{5 joint PhD publications} 
% $\rightarrow$ Left academia.

\item {Carlos Garcia Garcia} PhD at IFF-CSIC, 2017-20, co-supevisor.  %Defended 10-11-20
%\href{https://inspirehep.net/literature?sort=mostrecent&size=25&page=1&q=a%20c%20garcia-garcia%20and%20a%20zumalacarregui&ui-citation-summary=true}{4 joint PhD publications} 
% on cosmological tests of dark energy.
% Formal co-supervision with P. Ruiz-Lapuente
 % Balzan Fellow '19   %(2 Months visit Oxford).
 \\ $\longrightarrow$ \underline{Beecroft Fellow at U. Oxford (5y)}, Chair of Rubin/LSST's weak lensing \& LSS group.
 \\ $\longrightarrow$ \underline{Tenure track scientist} (\textit{Atraccion de Talento}) at CIEMAT
 
 \end{itemize}


\subsection*{PhD students: informal supervision}
\begin{itemize}\setlength{\itemsep}{1pt}
\item {Manish Tamta} (IIS Bangalore 2Mo stays in 2025, 206): lensing diffraction and dark matter.
  \item {Marienza Caldarola} (IFT-UAM/CSIC 3Mo stay in 2025): SBI lensing, 1 joint publication.
 \item {Hanxi Wang} (Oxford, 1.5Mo stay in 2024), lensing by binary black holes. 1 joint publication.
 \item {Martin Pijnenburg} (Geneva, 3Mo stay in 2024), collective diffraction effects.
\item {Luka Vujeva} (Copenhagen, 1 week visit in 2024), project on strong lensing. 1 joint publication.
 \item {Ho Yeuk Cheung} (JHU 2020-25, 1 week visit in 2023). \textbf{4 joint PhD publications (3 wo/ supervisor)} %lensing of gravitational waves by large-scale structure
 \\ $\to$    NASA-Hubble Fellowship (declined), \underline{Fellow at Princeton IAS} (w/ M. Zaldarriaga) 

 \item {Parita Mehta} (Warschau U. 1Mo stay in 2023), strong lensing of GWs
  \item {Dina Traykova} (2018-2021) (Oxford) %collaboration in dark energy projects \
  ({2 publications}) $\longrightarrow$ \underline{Postdoc at MPI-Grav}.
  \item Hector Villarrubia-Rojo (UCM) visited Berkeley 3Mo $\longrightarrow$ \underline{Postdoc at ETH} (w/ L Heisenberg).
    \item {Micah Brush} (UC Berkeley) %on tests of gravity with the CMB 
    ({1 publication}) $\longrightarrow$ \underline{Postdoc at U. Alberta}.
   \item {Alessio Spurio Mancini} (U. Heidelberg) %, on tests of gravity with weak gravitational lensing 
   ({1 publication})  $\longrightarrow$ \underline{Postdoc at U. Cambridge}.
   \item {Frank K\"onnig} and {Henrik Nersisyan} (U. Heidelberg) %, on massive gravity theories with ghosts 
   ({2 publications}).
%  \item \textbf{Manuel Trashorras} (PhD student at IFT-UAM) on modeling on cosmic voids and cross correlation of observables, started 2015 at IFT-UAM. 

 \item {Jose Maria Ezquiaga} (PhD at IFT-UAM/CSIC, 2015-19). \textbf{5 joint PhD publications (3 wo/ supervisor)}. %on gravitational waves and gravity theories  
% {Extraordinary Thesis Price} (UAM) and {Best Thesis} in Theoretical Physics (Spanish Physical Society) 2019 
\\
$\longrightarrow$ 2$\times$ thesis prizes, \underline{NASA Einstein Fellow at U. Chicago} (w/ D. Holz) \\ $\longrightarrow$ {Associate Prof. at NBI \& chair of LIGO-Virgo-KAGRA lensing group}. %24/383 = 0.062 chances of success
\end{itemize}




\subsection*{Master thesis supervision}
\begin{itemize}\setlength{\itemsep}{1pt}
  \item {Louis Salimbeni} (ETH Zurich, co-
supervision 2026)
  \item {Matteo Pagani} (U. Bologna, co-
supervision 2026)
  \item {Pablo Vega del Castillo} (LMU, supervision 2026)
  \item {Julius Streibert}  (FU Berlin, co-supervision 2022-23).  \textbf{1 publication} % on GW propagation beyond GR.
% $\longrightarrow$ \underline{PhD Student at Oxford U.}
\item {Shashwat Singh} (Paris Observatory, 2022-23, supervision). \textbf{1 publication} \\ $\longrightarrow$ \underline{PhD Student at Glasgow U.} %Project on probing ULDM with GWs.
  \item {Richard Stiskalek} (LMU, 2021-22, co-supervison). \textbf{2 publications} %on stong-field GW propagation in GR.
  \\ $\longrightarrow$ \underline{PhD at Oxford U.} (offers from U. Chicago, Monash U.), predoc \underline{Fellow at CCA} $\longrightarrow$ \underline{Postdoc at Oxford}

   \item {Alessandro Longo} (U. Trieste, 2020-21, co-supervision).  %Thesis on EFT-LSS in modified gravity. 
   \textbf{1 publication} \\
$\longrightarrow$ \underline{PhD Student at SISSA}

 \item {Janina Renk} (U. Heidelberg 2015-16, co-supervision). \textbf{2 publications} % on large-scale structure tests of gravity. % Project design and scientific co-supervision with L. Amendola. 
 \\ $\longrightarrow$ \underline{PhD at Stockholm U.} %w/ Joakim Edsj\"o. %on relativistic effects in alternative theories of gravity.
\end{itemize}

\subsection*{Master students: informal supervision}
\begin{itemize}\setlength{\itemsep}{1pt}
 \item {Selin Ustundag} (MPI-Grav 2026) internship on lensing by SMBHs.$\longrightarrow$ \underline{PhD at Newcastle}
  \item {Marcel Gl\"atzer} (HU Berlin, 2024-25).
  % \item {Vaibhav Omble} (LMU, 2023-25, co-supervison).
  \item {Jaime Raynaud} (TU Berlin \& KTH 2024-25) internship %on GW lensing. 
  $\longrightarrow$ \underline{PhD at MPIA} (w/ V. Springel)
  \item {Vaibhav Omble} (LMU 2023-24)
\end{itemize}

\subsection*{Undergraduate project supervision}

\begin{itemize}\setlength{\itemsep}{1pt}
\item  {Daryna Yushchenko} (Princeton Summer visitor's program, 2024)
 $\to$    \underline{PhD at Carnegie Mellon University}.
\item  {Lyla Choi} (Princeton Summer visitor's program, 2023) \textbf{1 publication} $\to$ \underline{PhD at Stanford}. %\underline{offers from Stanford, Chicago and Johns Hopkins U.}
 
\item  {Yukei Murakami} (UC Berkeley, 2018-19) %local vs cosmolgical evolution of scalar fields 
 $\to$    \underline{PhD at Johns Hopkins U.} (w/ Adam Riess - Nobel 2011).
\item {Ho Yeuk Cheung} (UC Berkeley, 2020) %lensing of gravitational waves by large-scale structure
 $\to$    \underline{PhD at Johns Hopkins University} (w/ Emanuele Berti).

    \item {Samuel Packman} (2021) (High-school student in Berkeley) $\to$ \underline{Undergraduate at MIT}.

%   \item \textbf{Antony Valeh} (UC Berkeley, 2018), %gravitational lensing of supernovae
%   \textbf{Giovanni Pecorino} (UC Berkeley 2018) % GW propagation beyond GR.
%   (other projects)
 \item {Brindha Kanniah} (MIT Summer visitors program to Heidelberg, 2015, supervision). %Project on gravitational waves in massive gravity theories 
 $\longrightarrow$ \underline{Masters at MIT}.
\end{itemize}

\subsection*{Other mentoring roles}\setlength{\itemsep}{1pt}
\begin{itemize}\setlength{\itemsep}{1pt}
 \item Mentor for the \textbf{LISA Early Career Scientists} (2021-25) %Kristen Lackeos, Jaime (check spelling)
\item \textbf{$\sim$ 80 Berkeley Connect undergraduate students} in 2 semesters (UC Berkeley, 2018): Group discussion and one-on-one mentoring.
\end{itemize}

\subsection*{Thesis panels}
  % \item {Caroline Jonas} (2020-2023) (PhD Student at MPI-Grav), member of \textbf{thesis advisory committee}.

\vspace{-.3cm}
\begin{longtable}{c p{15.5cm}}
2026 & Claire Rigouzzo's PhD defense (external examiner), Kings College (UK) \\
% 2025 & Hemantakumar Phurailatpam's PhD defense (external examiner), CUHK (Hong Kong) \\
2020-23 & Caroline Jonas' PhD thesis advisory comittee member, MPI-Grav, Potsdam (Germany) \\
2021 & {Vitor Manuel da Fonseca}'s Master's defense, University of Lisbon (Portugal) %Jan XX 2021
\\
 2020  & {Hector Villarubia-Rojo}'s PhD defense, Universidad Complutense de Madrid (Spain) %July 29, 2020
\end{longtable}
\mzc{Update!}


\subsection*{Scientific coordination}
% \vspace{-0.3cm}

\begin{itemize}

  \item \textbf{LISA Consortium} (2018--)
  \begin{itemize}
    \item 2025- Council member
    \item 2025- Senior Co-Chair, Cosmology Working Group.
    \item 2022 Section coordinator (``GW Lensing'') for LISA Cosmology White Paper (\href{https://inspirehep.net/literature/2066207}{link}).
  \end{itemize}

  \item \textbf{Euclid Consortium} (2014--)
  \begin{itemize}
    \item 2025- Lead for Science Working Group on Gravitational Waves (SWG-GWs)
    \item 2025- Lead for ``bright sirens'' (WP2) in Euclid's SWG-GW.
    \item Member of Strong Lensing, Cosmology Theory and Supernovae \& Transients SWGs.
  \end{itemize}
  \item \textbf{GW-Space 2050} defining ESA's Gravitational-Wave missions after LISA (2023-)
  \begin{itemize}
    \item Coordinator for ``Cosmography'' section.
  \end{itemize}
  \item \textbf{Multi-messenger lensing white paper} (2024) -- Coordinator for microlensing section (\href{https://inspirehep.net/literature/2904454}{link}).

  % \item GUEST

\end{itemize}

\subsection*{Event Organization}
\vspace{-.3cm}
\begin{longtable}{c p{14.5cm}}
2026
 & {``13th LISA Cosmology Working Group Workshop''} %at University of Barcelona, 
 Barcelona (Spain)  \href{https://indico.icc.ub.edu/event/677/}{link}. \\[3pt]
2025
 & {``12th LISA Cosmology Working Group Workshop''} %at Estonian Academy of Sciences, 
 Tallinn (Estonia)
 % : $\sim 60$  participants \& 25 recorded talks 
 \href{https://indico.cern.ch/event/1488828/}{link}. %\textbf{Funding: 10k\euro}
\\[3pt]
2024
 &  {``Lensing and Wave Optics in Strong Gravity Workshop''} at Erwin Schr\"odinger Institute, Vienna (Austria): $\sim 60$  participants \& 25 recorded talks \href{https://www.esi.ac.at/events/e546/}{link}. \textbf{Funding: 12k\euro}
\\[3pt]
2024 & {``LISA joint Fundamental Physics + Waveform meeting''} at Max Planck Institute for Gravitational Physics (MPI-Grav), Potsdam (Germany): 120+ participants, $\sim30$ talks \href{https://workshops.aei.mpg.de/fpmeetswavelisa/}{link} 
%\url{https://www.aei.mpg.de/139117/seminars} 
\\[2pt]
2020-22 &  {ACR Division Seminars} at Max Planck Institute for Gravitational Physics (MPI-Grav), Potsdam (Germany)  \href{https://www.aei.mpg.de/139117/seminars}{link} \\[2pt]
2016-17 & Weekly {BCCP Cosmology Lunch} at UC Berkeley (US) \\
2014-15 &  {Cosmology Seminars} at ITP, University of Heidelberg (Germany)  \href{http://www.thphys.uni-heidelberg.de/~cosmo/dokuwiki/doku.php/seminar}{link} \\[2pt]
2015  &  {Ninth TRR33 Winter School}, Passo de la Tonale (Italy) $\sim 40$ students (\href{http://darkuniverse.uni-hd.de/view/Main/WinterSchool15}{link}) \\[2pt]
&  {First Local Interactions Workshop} at ITP, Heidelberg (Germany) (\href{http://www.thphys.uni-heidelberg.de/~cosmo/inter15/}{link}) \\[2pt]
2014 &  {Eighth TRR33 Winter School}, Passo de la Tonale (Italy) $\sim 40$ students \href{http://darkuniverse.uni-hd.de/view/Main/WinterSchool14}{link} \\[2pt]
 & Organizer and lecturer {CLASS crash-course}, ITP Heidelberg \\[2pt]
2013 & {Seventh TRR33 Winter School}, Passo de la Tonale (Italy) $\sim 40$ students (\href{http://darkuniverse.uni-hd.de/view/Main/WinterSchool13}{link})
%	     &	with {David Mota}. Aug-Oct (two Months)
\end{longtable}


\section*{Teaching}
\subsection*{Graduate schools}
 \vspace{-.3cm}
 \begin{longtable}{c p{14.5cm}}
% AEI talks DB: 2024-09-09 | Gravitational Waves
  2024  & {Lecturer} at {TAE (Taller de Altas Energias)}, Benasque (Spain).
 \\ &
Course on Gravitational Waves (\href{https://www.benasque.org/2024tae/talks_contr/094_GW_Lectures_-_TAE_2024_-_Zumalacarregui_-_PDF.pdf}{slides}).
%  (\href{https://cescamilla4.wixsite.com/mctpcosmograv2}{website}).
\\[2pt]
 2018  & {Lecturer} at the {II Mesoamerican Workshop on Gravity and Cosmology}, Chiapas (Mexico). 
%  \\ & 
Course ``Tests of Gravity and Dark Energy with Cosmology and Gravitational Waves''
 (\href{https://cescamilla4.wixsite.com/mctpcosmograv2}{website}).
\\[2pt]
2017  & {Lecturer} at the {Cosmology School in the Canary Islands}, Fuerteventura (Spain).  Topic: ``alternative gravity theories'' (\href{http://iactalks.iac.es/m/talks/view/1081}{\underline{video link}},
 \href{http://www.iac.es/congreso/cosmo2017/}{website}) \\[2pt]
2017  & {Lecturer} at the {School for Cosmology Tools}, IFT Madrid (Spain). 
  % Lectures on CLASS, {\tt hi\_class} and Monte Python  \url{https://workshops.ift.uam-csic.es/cosmotools2017}
\\[2pt] 
%  & {Support Lecturer} (cancelled) at the {$\Lambda$CDM and Beyond}, Corfu (Greece) %\url{http://www.physics.ntua.gr/corfu2017/lc.html} %{$\Lambda$CDM and Beyond: Cosmology Tools in Theory and in Practice}, 
% \\[2pt] 
2014  & {Lecturer} at the {Special workshop on CLASS and MontePython}, MPIA Munich (Germany). %\\[-2pt] & 
% \url{https://lesgourg.github.io/courses.html}\\[2pt] 
\end{longtable}
%  \newpage

\noindent
\subsection*{University courses}
 \vspace{-.3cm}
\begin{longtable}{c p{14.5cm}}
2018-19 & {Berkeley Connect Instructor}:  mentoring and teaching (undergraduate level) at the University of California, Berkeley (2 semesters). %1-on-1 mentoring for 150+ undergraduates across 3 semesters, designed curricula, organized seminars, facilitated faculty panels, etc.
\\[2pt]
2014-15 & {Teaching assistant} for ``Theoretical Astrophysics'' (exercises, graduate level) course by M. Bartelmann at the University of Heidelberg. \\[2pt]
2014-15 & {Lecturer} in  ``Advanced Cosmology'', ``Introduction to CMB physics'', ``Statistical mechanics'' (graduate level) and ``advanced seminar'' (undergrad), University of Heidelberg. \\[2pt]
2013 & {Teaching assistant} for ``Gravitation and Cosmology'' and ``Physics of the Cosmos'' courses (undergraduate level) at the Universidad Autonoma de Madrid. \\[2pt]
2012& Mini-course on inhomogeneous cosmologies (undergraduate level) U. of Barcelona.\\
% 2001-2003& High school tutor in Physics, Mathematics and Chemistry.%(2 years).
\end{longtable}

\newpage

\subsection*{Training Received}
 \vspace{-8pt}
\begin{longtable}{c p{15.5cm}}
2025 & Basics of Project Management for Scientists %- MPI Training %- Successfully Master Your Project! -11.3. and 25.3. from 9:30 to 16:30 pm (12 hours)
\\[2pt]
2024 & {Start your Leadership Journey} - Max Planck Headquarters  %16 hours, 22.03.2024
% \\ & Max Planck Leadership - Part II
\\[2pt]
& Improved reading
\\[2pt]
2023 & {Recognizing Missconduct, Adddressing Misconduct} at MPI-Grav
\\[2pt]
2020 & {Unconscious Bias - Unconscious thought patterns} (by  Working Between Cultures) at MPI-Grav %4hours, 25.11.2020
\\[2pt]
% 2020 
& {Pathways to Scientific Teaching} (active, learner-centered instruction) at UC Berkeley %\\%[2pt]
% by Diane Ebert-May %Learner-centered instruction. 
\\[2pt]
2018 & {Ethics for Instructors} at UC Berkeley \\%[2pt]
 & Berkeley Connect training for Fellows at UC Berkeley. 
%  Introduction workshop plus modules on  \textit{Facilitating interactive discussions}, \textit{Storytelling \& science} 
\\[2pt]
% 2013 & Online course ``Accountable Talk: Conversation that Works'' on Coursera. 
\end{longtable}


\section*{Scientific presentations}

Over 130 scientific talks, including 13 plenaries at major conferences and 5 institute colloquia. 

\subsection*{Conferences \& workshops
 \hfill {\small  (invited$^\bigstar$, \href{https://miguelzumalacarregui.es/talks.html}{\underline{video link}}}) \quad} % \vspace{0.1cm} sub
\vspace{-.3cm}
\begin{longtable}{l p{7.cm} p{3.5cm} p{3cm}}
% AEI talks DB: 2026-07-13 | Gravitational Lensing of Waves: a new window into astrophysics, dark matter and gravity
2026  & Dark Matter \& Stars & Southampton, UK & \textbf{plenary talk$\bigstar$}\\
% AEI talks DB: 2026-04-22 | Novel tests of gravitational-wave propagation
& Dark Energy from Home & Online/Global &{\textbf{\href{https://youtu.be/1k0o0QiL288}{\underline{invited talk}}}} \\
% AEI talks DB: 2026-05-25 | Gravitational Lensing of Waves: a new window into astrophysics, dark matter and gravity
& EREP 2026 & Murcia, Spain & \textbf{plenary talk$\bigstar$}\\
% AEI talks DB: 2025-12-04 | Wave-optics phenomena, lensed gravitational waves and GW231123
2025 & GWPAW & Atlanta US &{contributed talk} \\
% AEI talks DB: 2025-07-16 | Wave-optics phenomena in lensed gravitational waves
& Amaldi/GR & Glasgow, Scotland &{contributed talk} \\
% AEI talks DB: 2025-06-24 | Discovering supermassive black-hole binaries through quasi-periodic lensing of starlight
   & Iberian Gravitational Waves Meeting & Madrid, Spain &{contributed talk} \\
% AEI talks DB: no matching entry
   & Space 2050 & Potsdam, Germany & contributed talk \\
%45 conferences up to Benasque 2024
% AEI talks DB: 2024-07-18 | Wave-optical Phenomena in GW Lensing
2024  & New Frontiers in Strong Gravity & Benasque, Spain & \textbf{{plenary talk}$^\bigstar$}
%2024 Pollica (declined)
\\
% AEI talks DB: 2024-06-15 | Lensing of Gravitational Waves New Opportunities for Fundamental Physics
 & Cosmology from Home & Online/Global & \href{https://youtu.be/cNc5VilkQH0?si=aw9emqQXW-r_2Nvr}{\underline{contributed talk}}
\\
2023  
% AEI talks DB: 2023-09-11 | On the propagation of gravitational waves: difraction, dispersion & birefringence [AEI location reads IFT, Madrid]
& COSMO '23 Conference & Madrid, Spain& contributed talk
\\
% AEI talks DB: no matching entry
& Very Light Dark Matter 2023 & Tokyo, Japan + online & 
\href{https://youtu.be/acbsylgDwQc?t=11484}{\underline{contributed talk}}
%https://indico.ipmu.jp/event/416/contributions/7487/
\\
% & Dark Side of the Universe & Rwanda  & plenary$^\bigstar$ (declined)
% \\
% & LISA cosmology & Stavangen, Norway  & Review$^\bigstar$ (declined)
% \\
% & EFTs meet gravitational waves & Benasque, Spain  & plenary$^\bigstar$ (declined)
% \\
2022  
% AEI talks DB: 2022-05-23 | Dark Energy [AEI title is just "Dark Energy"; the recording is titled "Dark energy and tests of gravity"]
& EuCAPT Annual Symposium & Online/Global  & \textbf{%
\href{https://cds.cern.ch/record/2810852}{\underline{plenary talk}}$^\bigstar$} 
\\
% AEI talks DB: 2022-08-30 | Gravitational wave lensing as a probe of cosmic structures
& EREP 2022 & Salamanca, Spain  & \textbf{{plenary talk}$^\bigstar$}
\\
% AEI talks DB: 2022-11-02 | GW lensing and fundamental physics
& VIII Meeting on Fundamental Cosmology & Granada, Spain  & \textbf{{plenary talk}$^\bigstar$} \\
% AEI talks DB: 2022-09-08 | Towards solutions to the Hubble problem beyond Einstein's Gravity
& Tensions in cosmology & Corfu, Greece  & \textbf{{plenary talk}$^\bigstar$} \\
% AEI talks DB: 2022-12-13 | On the propagation of gravitational waves: difraction, dispersion & birefringence
& Madrid Winter Workshop & Madrid, Spain / hybrid  & \textbf{{plenary talk}$^\bigstar$}%https://teorica.fis.ucm.es/madww2022/Presentation.html %Remote talk
% AEI talks DB: 2022-07-08 | Probing cosmic structures with gravitational wave lensing
\\ & $23^{\rm rd}$ Intl. Conference on GR \& Gravitation & Beijing/Online  & parallel talk
% AEI talks DB: 2022-07-09 | Probing cosmic structures with gravitational wave lensing [AEI location reads University of Glasgow, UK]
\\ & LISA Symposium & Online/Global  & \href{https://www.youtube.com/watch?v=XF_wsumMVKI}{\underline{parallel talk}}
\\
2021  
% AEI talks DB: 2021-04-17 | Gravitational Wave lensing beyond Einstein's General Relativity
& APS April Meeting
& Online/Global 
& parallel talk \\
% AEI talks DB: 2021-09-08 | Gravitational Wave lensing beyond Einstein's General Relativity
& Alternative Gravities \& Fundamental Cosmo. %sept 8 21
& Online/Global & contributed talk \\
% AEI talks DB: 2021-04-01 | Gravitational Wave lensing beyond Einstein's General Relativity [matched by elimination: the AEI database has three identical 2021 entries]
  & Iberian Cosmology Meeting (IberiCos)
& Online/Global 
& contributed talk \\
% AEI talks DB: 2020-12-10 | Tests of Gravity with Gravitational Waves [matched by elimination; the AEI entry is dated December, the CV files PONT under 2020]
2020  & Progress on Old and New Themes in Cosmo. %cosmology (PONT 2020)
& Online/Global %Avignon, France 
& \textbf{plenary talk$^\bigstar$} \\
% AEI talks DB: 2020-08-13 | Towards solutions to the Hubble problem beyond Einstein's Gravity
& Cosmology from Home & Online/Global  & \href{https://youtu.be/MHcLT31gaYQ}{\underline{parallel talk}} \\
% AEI talks DB: 2020-08-21 | GW lensing beyond General Relativity: Birefringence, Echoes and Shadows
& LISA Symposium & Online/Global  & \href{https://youtu.be/D3vSuuwIyfs}{\underline{parallel talk}} \\
2019  
& {COSMO '19 Conference} & Aachen, Germany &  \textbf{plenary talk$^\bigstar$}  \\
& {Lensing in the era of precision Cosmology} & BCCP, Berkeley, US &  contributed talk  \\
& {Theoretical Cosmology Meeting} & Leiden, NL & {invited talk$^\bigstar$} \\
2018  
& {Very High Energy Phenomena in the Universe} & Quy Nhon, Vietnam &  \textbf{plenary talk$^\bigstar$}  \\
& {Primordial vs Astrophysical origin of BHs} & CERN &  \href{https://cds.cern.ch/record/2320183}{\underline{talk + discussion}}$^\bigstar$  \\
& {Rencontres de Moriond} & La Thuile, Italy &  {contributed talk}  \\
& II Mesoamerican Workshop on Grav\&Cosmo % Gravity and Cosmology 
& MCTP, Mexico & \textbf{plenary lecturer$^\bigstar$} \\
      & {Direct Detection of Dark Energy} & Caltech, US & {contributed talk}  \\
      & {LSST DESC Meeting} & SLAC/Stanford, US & {short talk}  \\
      & {13$^{th}$ Iberian Cosmology meeting} & Lisbon, Portugal & {contributed talk}  \\
%  
2017   & {Dark Energy and Modified Gravity Cosmo.} & IPhT Saclay, France & \href{https://youtu.be/f47WCOd_Iqo}{\underline{contributed talk}}$^\bigstar$  \\
& {Theoretical Cosmology in the Light of Data} & Stockholm, Sweden & \href{https://youtu.be/zhuiV8sjFXs}{\underline{contributed talk}}  \\
2016 & {COSMO '16 Conference} & U. Michigan, US & parallel talk \\
& {Cosmology in the era of large scale surveys} & Florence, Italy & contributed talk \\
2015 & {Extended theories of gravity at Nordita} & Stockholm, Sweden & \href{https://www.youtube.com/watch?v=HmfRQPrbp3Y}{\underline{contributed talk}}$^\bigstar$ \\
  & {Cosmological tests of gravity workshop} & Leiden, Holland & contributed talk$^\bigstar$ \\
  & {Gravity at the largest scales} & Heidelberg, Germany & contributed talk \\
   & {Gravity Beyond Einstein} & Heidelberg, Germany & contributed talk \\
2014 & {Dark Energy Interactions at Nordita} & Stockholm, Sweden &\textbf{plenary talk$^\bigstar$} \\ 
 & {Benasque Modern Cosmology} & Benasque, Spain & contributed talk \\
 & {Essential Cosmology for the next generation} & San Lucas, Mexico & contributed talk
 \\[2pt] 
%  \arrayrulecolor{cyan}\hline
2013 & {Fundamental Physics in light of Planck \& DES} & Madrid, Spain & \textbf{plenary talk$^\bigstar$} \\
 & {8$^{th}$ Iberian Cosmology Meeting} & Granada, Spain & contributed talk \\
2012 & {Benasque Modern Cosmology} & Benasque, Spain & student talk \\
 & {7$^{th}$ Iberian Cosmology Meeting} & Lisbon, Portugal & contributed talk \\
2011 & {Spanish Relativity Meeting '11} & Madrid, Spain & contributed talk \\ 
 & {6$^{th}$ Iberian Cosmology Meeting} & Salamanca, Spain & contributed talk \\
2010 &  {5$^{th}$ Iberian Cosmology Meeting} & Porto, Portugal & contributed talk \\
\end{longtable}


\subsection*{Seminars \hfill {\small \sf  (invited$^\bigstar$, \href{https://miguelzumalacarregui.es/talks.html}{\underline{video link}})} } % \vspace{0.1cm} %Invited $\bullet$Invited 
\vspace{-.3cm}
\begin{longtable}{l p{7.cm} p{3.5cm} p{3cm}}
% AEI talks DB: 2026-04-21 | Across the Universe: GW231123 as a diffracted and magnified black hole merger
2026 & Princeton University &  Princeton, US  & IAS/PU GWs \\
% AEI talks DB: 2026-04-22 | Quasi-periodic lensing of starlight
& Institute for Advanced Study & Princeton, US & Astro coffee \\
% AEI talks DB: 2026-06-25 | Gravitational Lensing of Waves: a new window into astrophysics, dark matter and gravity
& Universidad Complutense - IPARCOS & Madrid, Spain & \textbf{\href{https://www.youtube.com/watch?v=VsK5Omc99DU}{\underline{Colloquium}}}$^\bigstar$ \\
% AEI talks DB: 2026-01-27 | A Magnified and Diffracted Black Hole Merger (GW231123)
 & Cosmology Talks (GW231123 lensing) & Online - Global & \href{https://youtu.be/ZGqOpDTL8k4?si=HyzJFQFJTZk-vvVy}{\underline{Recorded Webinar}} \\
% AEI talks DB: 2026-03-09 | Gravitational Lensing of Waves: a new window into astrophysics, dark matter and gravity
 & University of Virginia &  US (over zoom) & Gravity Seminar \\
% AEI talks DB: 2026-07-29 | Probing Low Mass Dark Matter Halos with LISA
 & Cosmology Talks (Stochastic diffraction) & Online - Global & \href{https://youtu.be/a5lIDTDRo4c}{\underline{Recorded Webinar}} \\

% AEI talks DB: 2026-05-22 | Gravitational Lensing of Waves: a new window into astrophysics, dark matter and gravity
 & University of Barcelona &  Spain & Cosmo Seminar \\
2025
% AEI talks DB: 2025-07-03 | Gravitational Lensing of Waves
& MPI-Grav. Albert Einstein Institute & Hannover, Germany & \textbf{Colloquium}$^\bigstar$ \\
% AEI talks DB: no matching entry
& IFT-UAM/CSIC & Madrid, Spain  & Institute Seminar \\ %& \textbf{Colloquium}$^\bigstar$ %or Institute seminar? \\
% \\
%74 seminars up to Copernicus!
% AEI talks DB: 2025-03-18 | Gravitational Lensing of Waves
& Copernicus Webinar & Online - Global & {Online Seminar}  \\
2024
% AEI talks DB: 2024-11-28 | Gravitational Lensing of Waves
& Leibniz Institute for Astrophysics & Potsdam, Germany & \textbf{Colloquium}$^\bigstar$  \\
% AEI talks DB: 2024-04-01 | Lensing of Gravitational Waves - New Opportunities for Physics
& Johns Hopkins University & Baltimore, US & Astro seminar \\
% AEI talks DB: 2024-02-28 | Lensing of Gravitational Waves New Opportunities for Fundamental Physics
& University of Birmingham & Birmingham, UK & Astro seminar$^\bigstar$ \\
% AEI talks DB: 2024-06-20 | Lensing of Gravitational Waves New Opportunities for Fundamental Physics
& Universidad Complutense & Madrid, Spain & Theory seminar$^\bigstar$ \\
% AEI talks DB: 2024-05-24 | Lensing of Gravitational Waves New Opportunities for Fundamental Physics
 & Donostia International Physics Center & San Sebastian, Spain & Online seminar$^\bigstar$ \\
% AEI talks DB: 2023-02-28 | On the propagation of gravitational waves: difraction, dispersion & birefringence
2023  & Oxford University & Oxford, UK & Cosmo Seminar$^\bigstar$ \\
% AEI talks DB: 2023-06-13 | On the propagation of gravitational waves: difraction, dispersion & birefringence
& Niels Bohr Institute & Copenhagen, DK & Gravity seminar$^\bigstar$ \\
% AEI talks DB: 2023-05-23 | On the propagation of gravitational waves: difraction, dispersion & birefringence
& Instituto de Fisica Corpuscular & Valencia, Spain & IFIC seminar$^\bigstar$ 
\\
% AEI talks DB: 2022-12-15 | On the propagation of gravitational waves: difraction, dispersion & birefringence
2022  & University of Vienna & Vienna, Austria & Physics seminar$^\bigstar$  %zoom
%\href{https://www.youtube.com/watch?v=w2vnenwbJWY}{\underline{Gravity Webinar}}%$^\bigstar$  %zoom
\\
% AEI talks DB: 2021-06-18 | The Dark Universe in the Era of Gravitational Wave Astronomy
2021 & {Institute for Theoretical Astrophysics}  & Oslo, Norway & \textbf{Colloquium}$\bigstar$ \\ 
% AEI talks DB: 2020-03-25 | Gravitational Wave lensing beyond Einstein's General Relativity [AEI date is WRONG: the talk was in 2021]
2021  & SISSA & Trieste, Italy & 
\href{https://www.youtube.com/watch?v=w2vnenwbJWY}{\underline{Gravity Webinar}}$^\bigstar$%
\\
% AEI talks DB: 2020-06-23 | Tests of gravity and dark energy with the propagation of gravitational waves [AEI date is WRONG: the talk was in 2021]
& University of Waterloo & Waterloo, Canada & 
Astro Seminar%zoom
\\
% AEI talks DB: 2020-04-01 | Dark Energy, the Hubble Problem and Gravitational-wave Propagation
2020  & Max Planck Institute for Gravitational Physics %(Albert Einstein Institute) 
& Potsdam, Germany & ACR seminar %zoom
\\
% AEI talks DB: 2020-04-15 | Dark Energy, the Hubble Problem and Gravitational-wave Propagatio [title truncated in the AEI database]
& Instituto de Astrof\'isica e Ci\^encias do Espa\c{c}o & Lisbon, Portugal & Cosmo Seminar %zoom
\\
% AEI talks DB: 2020-04-27 | Dark Energy, the Hubble Problem and Gravitational-wave Propagation
& ETH Zurich & Zurich, Switzerland & Cosmo Seminar %or Journal club %zoom
\\
% AEI talks DB: 2020-05-29 | Dark Energy, the Hubble Problem and Gravitational-wave Propagation
& University of Nottingham & Nottingham, UK & Cosmo Seminar %zoom
\\
2019  
& {University of California at Berkeley} & Berkeley CA, US & Astro Lunch (2$\times$) \\
& Lawrence Berkeley National Laboratory & Berkeley CA, US & DESI lunch (2$\times$) \\
& {Institut de Physique Th\'eorique, CEA} & Saclay, France & Particles \& Cosmo \\ 
& {D\'epartement d'Astrophysique, CEA} & Saclay, France & CosmoStat \\
& {Laboratoire de Physique Th\'eorique} & Orsay, France & Theory Seminar 
\\[2pt] 2018  
& SLAC National Accelerator Laboratory & Stanford U. CA, US & Theory seminar$^\bigstar$ \\ 
  & Max Planck Institute for Gravitational Physics 
%   (Albert Einstein Institute) 
  & Potsdam, Germany & ACR seminar$^\bigstar$ \\ 
& Berkeley Center for Theoretical Physics & Berkeley CA, US & 4D/Theory$^\bigstar$ \\
& University of Texas Austin & Austin TX, US & Theory seminar$^\bigstar$ \\
& Case Western Reserve University & Cleveland OH, US & Theory seminar$^\bigstar$ \\ %PACC
& University of Pittsburgh & Pittsburgh PA, US & Astro seminar$^\bigstar$ \\
  & University of California Santa Cruz & Santa Cruz CA, US & Cosmo Seminar$^\bigstar$ \\
   & University of Potsdam & Potsdam, Germany & Job talk$^\bigstar$ \\
  & {California Institute of Technology} & Pasadena CA, US & Cosmo Seminar \\ 
 & University of California Los Angeles & Los Angeles CA, US & Astro seminar \\
      & {University of Colorado at Boulder} & Boulder CO, US & Astro seminar \\ 
      & {Stanford University} & Palo Alto CA, US & KIPAC Tea talk \\ 
      & {ITP - University of Heidelberg} & Heidelberg, Germany & Cosmo Seminar$^\bigstar$ \\
   & {Institute for Theoretical Astrophysics} & Heidelberg, Germany & Blackboard Coll. \\
      & {University of California at Berkeley} & Berkeley CA, US & Astro Lunch (2$\times$) \\ 
    & {UCM - Universidad Complutense} & Madrid, Spain & Cosmo Seminar 
\\[2pt] 2017 &
{Institute for Applied Magnetism} & Madrid, Spain & \textbf{Colloquium}$^\bigstar$ \\
& Lawrence Berkeley National Laboratory & Berkeley CA, US & DESI lunch (2$\times$) \\
    & {Institute for Theoretical Physics UAM-CSIC} & Madrid Spain & gravity seminar \\
    & Instituto Avanzado de Cosmologia - {UNAM} & Mexico DF, Mexico & Cosmo Seminar \\
      & {Instituto de Estructura de la Materia - CSIC} & Madrid, Spain & Cosmo Seminar 
\\[2pt] 2016  & {KICP, University of Chicago} & Chicago, US & Cosmo Seminar$^\bigstar$ \\   
      & {CERN} & Geneva, Switzerland & Cosmo Seminar$^\bigstar$ \\ 
    & {Institute for Theoretical Physics UAM-CSIC} & Madrid Spain & Theory seminar \\
    & {University of Geneva} & Geneva, Switzerland & Cosmo Seminar \\
    & {Helsinki Institute of Physics} & Helsinki, Finland & Cosmo Seminar$^\bigstar$ \\
    & {Nordita/Oskar Klein Center} & Stockholm, Sweden & Theory seminar 
\\[2pt] 2015 & {VSI - University of Groningen} & Groningen, Holland & Theory seminar$^\bigstar$ \\  
& {Lawrence Berkeley National Laboratory} & Berkeley CA, US & INPA seminar \\
    & {Institute for Theoretical Astrophysics} & Heidelberg, Germany & Blackboard coll. \\
    & {ITP - University of Heidelberg} & Heidelberg, Germany & Cosmo Seminar \\
  & {ICC - University of Barcelona} & Barcelona, Spain & Cosmo Seminar \\
  & {UCM - Universidad Complutense} & Madrid, Spain & Cosmo Seminar 
  \\[2pt] 2014 
& {Perimeter Institute} & Waterloo, Canada & Cosmo Seminar$^\bigstar$ \\
& {Stanford University} & Palo Alto CA, US & Cosmo Seminar \\ 
  & {California Institute of Technology} & Pasadena CA, US & TAPIR seminar \\ 
  & {University of California Santa Cruz} & Santa Cruz CA, US & Cosmo Seminar \\ 
  & {University of California Irvine} & Irvine CA, US & Cosmo Seminar \\ 
  % & {Institute for Theoretical Astrophysics} & Heidelberg, Germany & blackboard colloquium \\ ??
  & {UCM - Universidad Complutense} & Madrid, Spain & Cosmo Seminar 
\\[2pt] 2013
& {University of Geneva} & Geneva, Switzerland & Cosmo Seminar$^\bigstar$ \\ 
 & {Lawrence Berkeley National Laboratory} & Berkeley CA, US & INPA seminar 
\\[2pt] 2012 
& {Helsinki Institute of Physics} & Helsinki, Finland & Cosmo Seminar$^\bigstar$ \\ 
& {Johns Hopkins University} & Baltimore MD, US & Whiteboard talk \\ 
  & {African Institute for Mathematical Studies} & Cape Town & Cosmo Seminar \\
  & {University of Cape Town} & South Africa & Cosmo Seminar \\  
  & {ITA - University of Oslo} & Oslo, Norway & Cosmo Seminar \\
  & {UCM - Universidad Complutense} & Madrid Spain & Cosmo Seminar \\ 
  & {Institute for Theoretical Physics UAM-CSIC} & Madrid Spain & Theory seminar
  \\[2pt] 2011 
  & {University of the Basque Country} & Bilbao, Spain & Cosmo Seminar$^\bigstar$ \\
  & {ITA - University of Oslo} & Oslo, Norway & Cosmo Seminar 
  \\[2pt] 2010  & {ITA University of Oslo} & Oslo, Norway & Cosmo Seminar \\
 & {ICC - University of Barcelona} & Barcelona Spain & Theory seminar \\
\end{longtable}


% \subsection*{\sc Other attended conferences, workshops and schools } % \vspace{0.1cm} $\bullet$
% \vspace{-8pt}
% \begin{longtable}{l p{10.5cm} p{5cm} }
% 2026 & New Horizons in Strong Field Gravity & Benasque, Spain\\
%  & Euclid meeting & Barcelon, Spain\\
% 2025 & Euclid \& ESLAB meeting & (remote participant) \\
% 2024 & LISA Fundamental Physics and Waveforms meeting & AEI Potsdam, Germany \\
% 2023 & Connecting the dots: Toward inspiral-merger-ringdown gravitational waveforms beyond general relativity & AEI Potsdam, Germany \\
% & LeadNet symposium for Max Planck Group Leaders & Berlin, Germany \\
% 2022 & {Cosmology from Home} & Online/Global \\ 
% 2021 & {Cosmology from Home} & Online/Global \\ 
% 2020 & {Probing Effective Ths of Gravity in Strong Fields and Cosmology} & KITP Santa Barbara \\ %https://www.kitp.ucsb.edu/activities/greft20
%  & {Rethinking the Relativistic Two-Body Problem} & AEI Potsdam, Germany \\ %https://workshops.aei.mpg.de/gwuniverse/
% 2019 & {Cosmology WG meeting, GdR Gravitational Waves} & APC Paris, France \\ %https://indico.in2p3.fr/event/19472
% 2018 & {Modelling the Extragalactic Sky} & UC Berkeley, US \\
% 2017 & {Cosmology with Neutral Hydrogen} & UC Berkeley, US \\
% %   & {Advances in Theoretical Cosmology in Light of Data} & Nordita, Sweden \\
% 2016 & {Unifying tests of General Relativity} & Caltech, US \\
%  & {Berkeley Neutrino Cosmology Workshop} & UC Berkeley, US \\
% & {Statistical sampling and non-sampling methods in cosmology} & UC Berkeley, US \\
%    & {Iberian Gravitational Waves meeting} & IFT Madrid, Spain \\
% 2015 & {Euclid Consortium Meeting} & EPFL Lausanne, Switzerland \\
%  2014 & {Structure Formation in the Modified Universe} & Leiden, Holland \\
% 2013 & {Cosmological Tests of Gravity} &  Oxford, UK. \\
%  &  {Cape Town International School of Cosmology} & Cape Town, South Africa \\
% 2011 & {The Dark Universe Conference} & Heidelberg, Germany \\ %Oct 
% 2010 & {Cosmo/CosPa} University of Tokyo & Tokyo, Japan \\  %Sept 
%  & {Benasque Modern Cosmology} &  Benasque, Spain \\
% 2009 & {Christmas Workshop'09} UAM & Madrid, Spain \\
%  & {4$^{th}$ Iberian Cosmology Meeting} & Madrid, Spain. \\
% 2008 & {Third Leopoldina Conference on Dark Energy} & Munich Germany \\
%  & {Second Universenet School and Meeting} & Oxford, UK \\
% 2007 & {Dark Energy Workshop} Dark Cosmology Center & Copenhagen, Denmark 
% \end{longtable}

\newpage

\section*{Outreach \& popular science}
\vspace{3pt}

% \textcolor{paragraphcolor}{\bf Talks \& Lectures}
% \vspace{-.2cm}
\subsection*{Talks \& Lectures}
\vspace{-.3cm}
\begin{longtable}{c p{15cm}}
% AEI talks DB: registered after the 2026-08-28 export, so absent from it
2026 & {Benasque Science Center} Talk: \textit{Los senderos tortuosos en la teoría de Einstein} (in Spanish, \href{https://campushuesca.unizar.es/noticia/charla-senderos-tortuosos-en-las-teorias-de-einstein}{link}) \\[2pt]
% AEI talks DB: 2025-09-13 | Sounds of silence: 10 years of gravitational-wave astronomy [AEI location reads Planetarium, Berlin]
2025 & {Long Night of Astronomy} Talk: \textit{Sounds of Silence: 10 years of GW astronomy} (\href{https://www.aei.mpg.de/1280508/10-years-of-gravitational-wave-astronomy}{link}) \\[2pt]
% AEI talks DB: 2025-05-10 | Mit Einstein auf krummen Wegen [AEI location reads AEI Potsdam]
& {Potsdam day of Science} Talk: \textit{Mit Einstein auf krummen Wegen} (in German, \href{http://ptdw.de/1037}{link}) \\[2pt]
% AEI talks DB: 2024-06-13 | Mit Einstein auf krummen Wegen [the AEI database has one entry for this and the talk below]
2024 & {Spendeprojekt}: presentation of the philanthropically supported project to donors \\
& Talk: \textit{Mit Einstein auf krummen Wegen} (in German) \\[2pt]
% AEI talks DB: 2024-11-09 | Deciphering Black Hole Symphonies
& Berlin Science Week event + Talk \textit{Deciphering Black Hole Symphonies} (\href{https://berlinscienceweek.com/programme/deciphering-black-hole-symphonies-new-world-gravitational-wave-astronomy}{link})
\\[2pt]
 2019 & {Astro Night}: monthly outreach event at the Astronomy Department, {UC Berkeley} CA, US. \\ & Talk: \textit{Gravitational Waves: Messengers from the Dark Universe} (\href{https://astro.berkeley.edu/astro-night}{link}).
\\[2pt]
 2018 & {Astrophysics Roundtable}: event for high-level donors of the Astronomy and Physics departments at {UC Berkeley} CA, US. Talk: \textit{Primordial black holes as a dark matter candidate}. % (\href{http://calday.berkeley.edu/detail.php?eventID=5011}{link})
\\[2pt]
 & {Cal Day 2018} at {UC Berkeley}, on the University's 150$^{\rm th}$ anniversary, Berkeley CA, US 
 \\ & \textit{Gravitational Waves: Messengers from the Dark Universe} (\href{http://calday.berkeley.edu/detail.php?eventID=5011}{link})
\\[2pt]
& {Taste of Science Festival}, Berkeley %CA, US 
% \\ & 
\textit{Crossing the Event Horizon:
Mysteries of Gravity} (\href{https://tasteofscience.org/san-francisco-events/2018/4/24/going-down-diving-into-gravity-and-the-deep-blue-sea}{link})
% \\[2pt]
\end{longtable}

% \newpage
% \noindent
% \textcolor{paragraphcolor}{\bf Articles (In Print)}
% \vspace{-0.2cm}
\subsection*{Articles (In Print)}
\vspace{-.3cm}
\begin{longtable}{c p{15cm}}
2018  & {Article} on {Investigacion y Ciencia} (Spanish edition of Scientific American) \href{https://www.investigacionyciencia.es/revistas/investigacion-y-ciencia/un-nuevo-plancton-737/las-teoras-de-la-gravedad-tras-la-tormenta-csmica-16414}{\it Las teor\'ias de la gravedad tras la tormenta c\'osmica} (in Spanish).
\\[2pt]
  & {Article} in {Berkeley Science Review} \href{https://www.berkeleysciencereview.com/article/2018/05/19/from-the-field}{\it From the Field} (Spring 2018 Issue).
\\[2pt]
  & {2 Articles} in the {Mathematical \& Physical Sciences Newsletter}, {College of Letters \& Sciences}, UC Berkeley.
 \end{longtable}
\vspace{-0.1cm}

% \noindent
% \textcolor{paragraphcolor}{\bf Articles (Online)}
% \vspace{-0.2cm}

\subsection*{Articles (Online)}
\vspace{-.3cm}
\begin{longtable}{c p{15cm}}
2021 & {Article} \textit{``Get your application noticed''} in \href{https://www.berkeleysciencereview.com/article/2021/05/22/get-your-application-noticed}{Berkeley Science Review Blog} %  
\\[2pt]
2018 & {Article} \textit{``Ondas gravitacionales y energ\'ia oscura: el fin de una era''} in \href{https://lafisicadelgrel.blogspot.com/2018/06/ondas-gravitacionales-y-energia-oscura.html}{``La F\'isica del Grel'' Blog} %  
(in Spanish)  \\[2pt]
2016
& {Article} {\it ``Las ondas gravitacionales: siguiendo el rumor mas tenue hacia los secretos del cosmos''} FronteraD
(in Spanish, \href{https://www.fronterad.com/las-ondas-gravitacionales-siguiendo-el-rumor-mas-tenue-hacia-los-secretos-del-cosmos/}{link})
\\[2pt]
& {\it ``GW150914, la persistencia y el genio''} (in Spanish, \href{https://www.fronterad.com/gw150914-la-persistencia-y-el-genio/}{link})\\[2pt]
& Comment on GW detection featured in ``Divulgacion cientifica de cientificos'' (in Spanish, \href{https://divulgacioncientificadecientificos.blogspot.com/2018/10/ondas-gravitacionales-miguel.html}{link})
\\[2pt]
2015
 &  {Article} {\it ``El mundo en escala logar\'itmica''}  FronteraD
(in Spanish,
\href{https://www.fronterad.com/pequeno-mapa-del-mundo-en-potencias-de-diez-i-contando-cifras/}{link}) %in Oct. 2015 
\\[2pt]
 &  {\it ``El valor de la ficci\'on''}  FronteraD
(in Spanish,
\href{https://www.fronterad.com/el-valor-de-la-ficcion/}{link}) %in Oct. 2015 
\\[2pt]
 &  {\it ``Persistencia y Serendipia''}  FronteraD
(in Spanish,
\href{https://www.fronterad.com/persistencia-y-serendipia/}{link}) %in Oct. 2015 
\\[2pt]
2011 & {Article} {\it ``Los neutrinos y el l\'imite de velocidad de Einstein''} % (article in Spanish)
%\\ & 
FronteraD (in Spanish, \href{https://www.fronterad.com/los-neutrinos-y-el-limite-de-velocidad-de-einstein/}{link})  and El Faro (\href{https://www.elfaro.net/es/201110/el_agora/6219/}{link}) 
% \\[2pt]
%  - & Outreach {talks for high-school students} at the Park School of Baltimore (US) and Colegio Montserrat (Madrid, Spain)
\end{longtable}

%TODO: rethink categories
% \subsection*{\sc Articles Authored (Online)}
% \vspace{-0.3cm}
% \begin{longtable}{c p{15cm}}
%  
% \end{longtable}

% \newpage

% \noindent
% \textcolor{paragraphcolor}{\bf Radio/Podcasts}
% \vspace{-0.2cm}

\subsection*{Radio \& podcasts}
\vspace{-.3cm}
\begin{longtable}{c p{15cm}}
 2023 & {Presentation on lensing of continuous gravitational waves} with Jose Edelstein %in \href{https://youtu.be/lyETW384gE8}{{\it Amautas} (radio)} 
%  (in Spanish).
\\[2pt]
 2018 & {Interview on dark matter, black holes and gravitational lensing of supernovae} with Jaime Garcia in \href{https://youtu.be/lyETW384gE8}{{\it Segmento Astronomico} (radio)} 
%  (in Spanish).
\\[2pt]
 & {Program on axial-gravitational anomaly} in \href{https://www.ivoox.com/podcast-horizonte-sucesos-cine-ciencia-entre_sq_f1432346_1.html}{{\it Horizonte de Sucesos} (podcast)} 
%  (in Spanish).
\\[2pt]
2017  & {Interview about Gravitational Waves} in \href{https://www.spreaker.com/user/horizontedesucesos/horizonte-de-sucesos-ondas-gravitacional}{{\it Horizonte de Sucesos} (podcast)} (in Spanish).
\\[2pt]
 & {Interview about Dark Energy} in \href{https://www.spreaker.com/user/horizontedesucesos/la-energia-oscura-2-parte-14-09-17}{{\it Horizonte de Sucesos} (podcast)} (in Spanish).
\\[2pt]
& Research featured in the {Coffee Break podcast} (in Spanish): \href{https://xn--sealyruido-u9a.com/?p=1216}{Gravitational waves and dark energy} (minute 47) and \href{https://xn--sealyruido-u9a.com/?p=1285}{Primordial Black Holes} (minute 72).  \\[2pt]
\end{longtable}
\vspace{-0.3cm}

% \newpage
% \noindent
% \textcolor{paragraphcolor}{\bf Media Features and Interviews}
% \vspace{-0.2cm}

\subsection*{Media Features and Interviews}
\vspace{-.3cm}
\begin{longtable}{c p{15cm}}
2026 & Feature in {\it Science} \href{https://www.science.org/content/article/extreme-black-hole-merger-may-be-cosmic-illusion}{\it Extreme black hole merger may be a cosmic illusion}, by Daniel Clery \\[2pt]
& Feature in {\it bild der wissenschaft} \href{https://wissenschaft.de/artikel/raetsel-um-uebergewichtige-schwarze-loecher-geloest}{\it Rätsel um ``übergewichtige'' Schwarze Löcher gelöst?}, by Nadja Podbregar (in German) \\[2pt]
& Feature in {\it National Geographic España} \href{https://www.nationalgeographic.com.es/ciencia/agujeros-negros-mas-masivos-jamas-vistos-podrian-estar-enganandonos-lente-cosmica-habria-hecho-que-parecieran-gigantes_29335}{\it Los agujeros negros más masivos jamás detectados podrían ser una ilusión causada por una lente gravitacional}, by Sarah Romero (in Spanish) \\[2pt]
& Research featured, by language, in \href{https://phys.org/news/2026-09-biggest-black-hole-merger-massive.html}{Phys.org} \& \href{https://skycr.org/2026/09/03/gw231123-lensing-black-hole-merger/}{SKYCR} (English), \href{https://www.scinexx.de/news/kosmos/ist-das-geheimnis-uebergewichtiger-schwarzer-loecher-gelueftet/}{scinexx} (German), \href{https://www.xataka.com/espacio/fusion-agujeros-negros-que-rompio-todas-reglas-tenia-truco-lente-gravitacional-pudo-haber-inflado-su-tamano}{Xataka} (Spanish), \href{https://trustmyscience.com/les-trous-noirs-interdits-seraient-ils-une-illusion-une-etude-propose-une-explication-surprenante/}{Trust My Science} (French), \href{https://www.tomshw.it/scienze/una-fusione-record-di-buchi-neri-potrebbe-aver-ingannato-ligo}{Tom's Hardware} (Italian), \href{https://www.ccvalg.pt/astronomia/noticias/2026/09/1_gw231123.htm}{Centro Ci\^encia Viva} (Portuguese), \href{https://wszystkoconajwazniejsze.pl/pepites/najwieksze-zderzenie-czarnych-dziur-moglo-nie-byc-rekordowe-sygnal-mogl-zostac-zaklocony/}{Wszystko co Najwa\.zniejsze} (Polish) \& \href{https://srpske.rs/en/news/nauka/2026/09/05/gw231123-record-black-holes-may-be-lensing-illusion}{Srpske} (Serbian). \\[2pt]
% Coverage of the same paper that is not listed above, so the section stays
% one screen. Restore any of these into the row if a call asks for reach:
%   Nueva Pensamiento Critico -- a full republication of the National Geographic
%     piece, not independent coverage:
%     https://nuevapensamientocritico.org/2026/09/06/los-agujeros-negros-mas-masivos-jamas-detectados-podrian-ser-una-ilusion-causada-por-una-lente-gravitacional/
%   Xataka also cut a video short: https://www.youtube.com/shorts/6lxAV2ypYZ0
%   National Geographic Espana's own post:
%     https://www.facebook.com/NationalGeographicEsp/posts/1468589031969000/
%   The shortlink https://dozz.es/v7qck resolves to the National Geographic
%     article linked above; the canonical URL is used instead, since a
%     shortener outlives nothing.
& Press release {\it Do black holes pretend to have large masses?} \href{https://www.aei.mpg.de/1490509/do-black-holes-pretend-to-have-large-masses}{MPI-Grav} (also in \href{https://www.aei.mpg.de/1490708/do-black-holes-pretend-to-have-large-masses}{German}) \& \href{https://www.mpg.de/26940743/do-black-holes-pretend-to-have-large-masses}{Max Planck Society} \\[2pt]
& Press release {\it New method could reveal hidden supermassive black hole binaries} \href{https://www.aei.mpg.de/1407139/new-method-could-reveal-hidden-supermassive-black-hole-binaries}{MPI-Grav} (also in \href{https://www.aei.mpg.de/1407247/new-method-could-reveal-hidden-supermassive-black-hole-binaries}{German}) \& \href{https://www.physics.ox.ac.uk/news/new-method-could-reveal-hidden-supermassive-black-hole-binaries}{University of Oxford} \\[2pt]
& Research featured in \href{https://phys.org/news/2026-02-gravitational-lensing-technique-unveils-supermassive.html}{Phys.org}, \href{https://www.thebrighterside.news/post/physicists-propose-a-new-way-to-spot-supermassive-black-hole-pairs/}{Bright Side News}, \href{https://www.myscience.org/en/news/2026/new_method_could_reveal_hidden_supermassive_black_hole_binaries-2026-mpg}{myScience}, \& \href{https://scienmag.com/new-technique-may-uncover-hidden-supermassive-black-hole-pairs/}{ScienMag}. \\[2pt]
& Podcast \& article \href{https://www.rexmolon.es/rexmolon-producciones-noticias-de-generacion-de-contenidos-de-divulgacion-cientifica-y-tecnologica-y-de-innovacion-y-desarrollo/un-nuevo-mtodo-permite-descubrir-agujeros-negros-gemelos-escondidos-en-el-corazn-de-las-galaxias}{\it Un nuevo método permite descubrir agujeros negros gemelos escondidos en el corazón de las galaxias} (in Spanish)\\[2pt]
& Feature in the \href{https://www.lisamission.org/?na=v&nk=882-dcfc753f05&id=31}{LISA Newsletter} \\[2pt]
2025  & Press release MPI-Grav 
\href{https://www.aei.mpg.de/1359867/erc-consolidator-grant-for-miguel-zumalac-rregui}{\it Around 2.2 million euros to open a new window on the Universe} (also in \href{https://www.aei.mpg.de/1363635/erc-consolidator-grant-for-miguel-zumalac-rregui}{German})
\\[2pt]
& Press release Max Planck Society
\href{https://www.mpg.de/25845263/erc-consolidator-grants-2025}{\it Six ERC Consolidator Grants for Max Planck} (also in \href{https://www.mpg.de/25844813/sechs-erc-consolidator-grants-fur-max-planck}{German})
\\[2pt]
 & Article in El Pa\'is \href{https://elpais.com/ideas/2025-08-02/velocidades-improbables-y-agujeros-de-gusano-podremos-viajar-en-el-tiempo.html}{\it Podremos viajar en el tiempo?} (in Spanish)
\\[2pt]
& Press release MPI-Grav \href{https://www.aei.mpg.de/1227766/with-einstein-on-crooked-paths}{\it With Einstein on crooked paths} (also in \href{https://www.aei.mpg.de/1227958/with-einstein-on-crooked-paths}{German})\\[2pt]
& Press release MPI-Grav \href{https://www.aei.mpg.de/1240576/euclid-opens-data-treasure-trove-offers-glimpse-of-deep-fields}{\it Euclid opens data treasure trove, offers glimpse of deep fields} (also in \href{https://www.aei.mpg.de/1240587/euclid-opens-data-treasure-trove-offers-glimpse-of-deep-fields}{German})\\[2pt]
& Kurzgesagt video \href{https://www.youtube.com/watch?v=zozEm4f_dlw}{Cosmic Crisis} -- Expert advisor (\href{https://sites.google.com/view/sources-cosmic-crisis/}{link}) \\[2pt]
& Swiss TV presentation by collaborator  \href{https://www.unige.ch/cite/evenements/ma-these-en-180-secondes/mt180-2025/les-finalistes/martin-pijnenburg}{\it Au travers des lunettes de l'Univers} (in French) \\[2pt]
2024 & Max Planck Society Yearbook (2023) \href{https://www.mpg.de/21438017/aei_jb_2023}{\textit{Looking at gravitational waves through Einstein's lenses}} \\[2pt]
& Research highlight MPI-Grav \href{https://www.aei.mpg.de/1183585/strong-space-time-curvature}{\textit{How does strong space-time curvature affect the propagation of GWs?}} \\[3pt]
2023 & {Astrobites feature} \href{https://astrobites.org/2023/04/11/gravitational-waves-a-la-general-relativity-or-scrambled/}{\textit{gravitational waves \`{a} la general-relativity or scrambled}} \\[2pt]
& Research featured in video \href{https://www.youtube.com/watch?v=H-nGZCdNuRQ}{\textit{El Nombre Mas Comico de un Articulo Cientifico}} (2.2M views)  \\[2pt]
2022 & {Community Profile} in the \href{https://www.eucapt.org/post/community-profile-miguel-zumalacarregui}{EuCAPT Blog} \\[2pt]
2020 & {Interview for El Confidencial} with Hector Barn\'es \href{https://www.elconfidencial.com/tecnologia/ciencia/2020-01-18/no-ligo-macho-articulo-cientifico-zumalacarregui_2411428/}{\textit{No LIGO MACHO: el t\'itulo m\'as cachondo de un art\'iculo cient\'ifico es idea de un espa\~nol}} (in Spanish).
\\[2pt]
 & {Interview for Forbes} with Ethan Siegel \href{https://www.forbes.com/sites/startswithabang/2020/12/24/how-gravitational-waves-might-wind-up-proving-einstein-wrong/?sh=149999106845}{\textit{How Gravitational Waves Might Wind Up Proving Einstein Wrong}}.
 \\[2pt]
 & {Research highlight MPI-Grav} \href{https://www.aei.mpg.de/597240/testing-einstein-s-theory-with-lensing-of-gravitational-waves}{\textit{Testing Einstein's theory with lensing of gravitational waves}}.
\\[2pt]
 & {Press Release at University of Chicago} \href{https://news.uchicago.edu/story/ripples-space-time-could-provide-clues-missing-components-universe}{\textit{Ripples in space-time could provide clues to missing components of the universe}}. 
\\[2pt]
 & {Research Featured} in \href{https://thedebrief.org/gravitational-waves-offer-clues-to-a-longstanding-cosmic-mystery/}{\textit{The Debrief}}
 and \href{https://www.europapress.es/ciencia/astronomia/noticia-encontrar-piezas-faltan-comprension-universo-20201228165225.html}{\textit{Europa Press}} (in Spanish). 
\\[2pt]
2018 & {Interview for Wired \& Quanta Magazine} with Katia Moskvitch, quoted multiple times on \href{https://www.quantamagazine.org/troubled-times-for-alternatives-to-einsteins-theory-of-gravity-20180430/}{\textit{Troubled Times for Alternatives to Einstein's Theory of Gravity}} (Published also in \href{https://www.wired.com/story/troubled-times-for-alternatives-to-einsteins-theory-of-gravity/}{Wired}).
\\[2pt]
& {Interview for UC Berkeley News} with Robert Sanders \href{http://news.berkeley.edu/2018/10/02/black-holes-ruled-out-as-universes-missing-dark-matter/}{\textit{Black holes ruled out as universe's missing dark matter}} (also on \href{http://physics.berkeley.edu/news-events/news/20181002/black-holes-ruled-out-as-universes-missing-dark-matter}{Berkeley Physics}).
\\[2pt]
& {Interview for Wired} with Sophia Chen, quoted on \href{https://www.wired.com/story/for-dark-matter-hunters-out-there-theories-are-catching-on/}{\textit{Dark Matter Hunters Pivot After years of Failed Searches}}.
\\[2pt]
& {Interview for The Daily Californian} with Yao Huang \href{http://www.dailycal.org/2018/10/03/uc-berkeley-researchers-find-black-holes-do-not-constitute-all-dark-matter/}{\textit{UC Berkeley researchers find black holes do not constitute all dark matter}}.
\\[2pt]
& {Interview for PNAS feature} with Anil Ananthaswamy, \href{http://www.pnas.org/content/115/40/9810.full}{\textit{
Inner Workings: How fast is the universe expanding? Clashing measurements may point to new physics}}.
\\[2pt]
& {Press releases at CEA} (in French) \href{http://www.cea.fr/drf/Pages/Actualites/En-direct-des-labos/2018/matiere-noire--la-theorie-des-trous-noirs-disqualifiee-.aspx}{\textit{Mati\`ere noire: la th\'eorie des trous noirs disqualifi\'ee}} \& \href{https://www.ipht.fr/Phocea/Vie_des_labos/News/index.php?id_news=860}{\textit{Les exp\'eriences LIGO/VIRGO contraignent-elles l'\'energie noire?}}.
\\[2pt]
& Mention in \href{https://mathscholar.org/2018/04/will-physics-anomalies-lead-to-new-physics/}{\it Will physics anomalies lead to new physics?} in Mathscholar
\\[2pt]
2017 
& {Interview for Science News} with Lisa Grossman, quoted on \href{https://www.sciencenews.org/article/what-detecting-gravitational-waves-means-expansion-universe}{\textit{What detecting Gravitational waves means for the expansion of the universe}}.
\\[2pt]
& {Interview for New Scientist} with Anil Ananthaswamy, quoted on \href{https://www.newscientist.com/article/2151020-dark-energy-survives-neutron-star-crash-test-while-rivals-fail/}{\textit{Dark energy survives neutron star crash test while rivals fail}}.
\\[2pt]
& {Interview for Lawrence Berkeley National Laboratory News} with Glenn Roberts Jr. \href{http://newscenter.lbl.gov/2017/12/18/star-mergers-test-gravity-dark-energy-theories/}{Star Mergers: A New Test of Gravity, Dark Energy Theories}.
\\[2pt]
& {Interview for Nature} with Davide Castelvecchi (unpublished). %quoted on \href{}{\textit{xx}}
\\[2pt]
& {Research featured in New Scientist}, \href{https://www.newscientist.com/article/2156252-ancient-black-holes-ruled-out-as-source-of-all-dark-matter}{\textit{Ancient black holes ruled out as source of all dark matter}}.
\\[2pt]
& Research featured in {Forbes} in \href{https://www.forbes.com/sites/startswithabang/2017/10/25/merging-neutron-stars-deliver-deathblow-to-dark-matter-and-dark-energy-alternatives}{\textit{Merging Neutron Stars Deliver Deathblow To Dark Matter And Dark Energy Alternatives}} and
\href{https://www.forbes.com/sites/startswithabang/2017/12/12/dark-matter-winners-and-losers-in-the-aftermath-of-ligo}{\textit{Dark Matter Winners and Loosers in the Aftermath of LIGO}}.
\\[2pt] 
& {Interview} for \href{http://www.rewisor.com/ondas-gravitacionales-estrellas-neutrones}{\texttt{Rewisor} on GW170817} (in Spanish)  \\[2pt]
 &  Profile featured in the \href{https://medium.com/sparrho/young-researchers-present-on-the-global-stage-part-7-8a31da24c8ab}{Sparrho blog} as finalist for the Early Career Researcher Prize
\\[2pt] 
& Researach features: paper on gravitational waves in \href{http://francis.naukas.com/2017/12/20/la-velocidad-de-las-ondas-gravitacionales/#comment-445151}{Naukas}, article on primordial black holes in \href{http://francis.naukas.com/2017/12/07/42781/}{Naukas} \& \href{http://curioseantes.blogspot.com.es/2017/12/hay-suficientes-macho-o-no.html?m=1}{Curioseantes} \\[2pt]
2016 
& {2015's 5th best Swedish research} in space sciences: featured scientific publication (\href{http://arxiv.org/abs/1502.02666}{Bettoni \& MZ PRD '15}) by \href{http://www.popularastronomi.se/2016/01/arets-basta-svenska-astroforskning-2015-mork-materia-utmanare-till-einstein-och-rosettas-komet/}{\tt www.popularastronomi.se} (in Swedish)
%(Bettoni \& Zumalacarregui PRD '15). %Bettoni \& Zumalacarregui (PRD '15) 1505.03839
\\[2pt]
 & Popular science recommendations featured in \href{http://divulgacioncientificadecientificos.blogspot.com.es/2016/03/miguel-zumalacarregui-perez-nordita.html}{\tt divulgacioncientificadecientificos}
\\[2pt]
2013 & PhD thesis featured on the science blog \href{http://francis.naukas.com/2013/01/01/el-metodo-cientifico-la-energia-oscura-y-la-importancia-de-la-quintaesencia/}{\tt www.francis.naukas.com} \\[2pt]
\end{longtable}
\vspace{-0.3cm}


% \section*{Scientific publications}\label{sec:publication_list}

% 58 items w/ 7,174 citations (\href{https://inspirehep.net/authors/1113379?ui-citation-summary=true}{InSPIRE}). {\signcited{} Top 1\% in Physics,  \signhot{} Top 0.1\% \& $\leq 2$ years {(\href{https://www.webofscience.com/wos/woscc/summary/6ce91b13-02c0-414c-ac15-c4f89814dc78-fc3901c0/relevance/1}{Web of Science})}}

% \subsection*{Peer reviewed articles}

% \begin{itemize}\setlength{\itemsep}{1pt}

% \item[\Num] \textit{``Reconstructing the dark energy density in light of DESI BAO observations''} M.~Berti, E.~Bellini, C.~Bonvin, M.~Kunz, M.~Viel, M.~Zumalacarregui \href{https://arxiv.org/abs/2503.13198}{2503.13198}. 
% \\
% PRD accepted \textbf{Editor's suggestion}
% \hfill \underline{16 citations}

% \item[\Num] \textit{``Gravitational wave lensing: probing Fuzzy Dark Matter with LISA,''}  
% S.~Singh, G.~Brando, S.~Savastano and M.~Zumalacarregui,  
% \href{http://arxiv.org/abs/arXiv:2502.10758}{arXiv:2502.10758}.  JCAP accepted 
% \hfill \underline{3 citations}

% \item[\Num] \textit{``Gravitational wave propagation beyond General Relativity: geometric optic expansion and lens-induced dispersion,''}  
% N.~Menadeo and M.~Zumalacarregui,  \\ {PRD (accepted)}
% \href{http://arxiv.org/abs/arXiv:2411.07164}{arXiv:2411.07164}.  
%  \hfill  \underline{1 citation} 

% \item[\Num] \textit{``GLoW: novel methods for wave-optics phenomena in gravitational lensing,''}  
% H.~Villarrubia-Rojo, S.~Savastano, M.~Zumalacarregui L.~Choi, S.~Goyal, L.~Dai and G.~Tambalo, \\ {PRD (accepted)}
% \href{http://arxiv.org/abs/arXiv:2409.04606}{arXiv:2409.04606}.  
%  \hfill  \underline{15 citations} 

% %\cite{Brando:2024inp}
% \item[\Num] \textit{``Signatures of dark and baryonic structures on weakly lensed gravitational waves,''}  
% G.~Brando, S.~Goyal, S.~Savastano, H.~Villarrubia-Rojo and M.~Zumalac\'arregui,  
% Phys. Rev. D \textbf{111}, no.2, 024068 (2025)  
% % doi:10.1103/PhysRevD.111.024068  
% \href{http://arxiv.org/abs/arXiv:2407.04052}{arXiv:2407.04052}  
% \hfill  \underline{5 citations}
% %5 citations counted in INSPIRE as of 24 Feb 2025


% \item[\Num] \textit{``Gravitational-wave lensing in Einstein-aether theory,''}
% %\cite{Streibert:2024cuf}
% J.~Streibert, H.~O.~Silva and M.~Zumalacarregui,
% \href{http://arxiv.org/abs/arXiv:2404.07782}{arXiv:2404.07782}. PRD accepted. \hfill \underline{4 citations}

% %\cite{Cheung:2024ugg}
% \item[\Num] \textit{``Probing minihalo lenses with diffracted gravitational waves,''}
% M.~H.~Y.~Cheung, K.~K.~Y.~Ng, M.~Zumalac\'arregui and E.~Berti,
% Phys. Rev. D \textbf{109}, no.12, 124020 (2024)
% % doi:10.1103/PhysRevD.109.124020
% \href{http://arxiv.org/abs/arXiv:2403.13876}{arXiv:2403.13876}
% \hfill  \underline{10 citations}
% %0 citations counted in INSPIRE as of 22 Jun 2024


% %\cite{Oancea:2023hgu}
% \item[\Num] \textit{``Probing general relativistic spin-orbit coupling with gravitational waves from hierarchical triple systems,''}
% M.~A.~Oancea, R.~Stiskalek and M.~Zumalac\'arregui. \textbf{MNRAS Letters} 535 (2024) \\
% \href{http://arxiv.org/abs/arXiv:2307.01903}{arXiv:2307.01903}.
% \hfill  \underline{17 citations}
% %9 citations counted in INSPIRE as of 31 May 2024


% %\cite{Savastano:2023spl}
% \item[\Num] \textit{``Weakly lensed gravitational waves: Probing cosmic structures with wave-optics features,''}
% S.~Savastano, G.~Tambalo, H.~Villarrubia-Rojo and M.~Zumalacarregui,
% Phys. Rev. D \textbf{108}, no.10, 103532 (2023)
% % doi:10.1103/PhysRevD.108.103532
% \href{http://arxiv.org/abs/arXiv:2306.05282}{arXiv:2306.05282}
% %15 citations counted in INSPIRE as of 25 Jun 2024
% \hfill  \underline{25 citations}



% %\cite{Goyal:2023uvm}
% \item[\Num] \textit{``Probing lens-induced gravitational-wave birefringence as a test of general relativity,''}
% S.~Goyal, A.~Vijaykumar, J.~M.~Ezquiaga and M.~Zumalacarregui,
% Phys. Rev. D \textbf{108}, no.2, 024052 (2023)
% % doi:10.1103/PhysRevD.108.024052
% \href{http://arxiv.org/abs/arXiv:2301.04826}{arXiv:2301.04826}
% %16 citations counted in INSPIRE as of 26 Jun 2024
% \hfill  \underline{26 citations}

% %\cite{Savastano:2022jjv}
% \item[\Num] \textit{``Through the lens of Sgr A*: Identifying and resolving strongly lensed continuous gravitational waves beyond the Einstein radius,''}
% S.~Savastano, F.~Vernizzi and M.~Zumalac\'arregui, \\
% Phys. Rev. D \textbf{109}, no.2, 024064 (2024)
% % doi:10.1103/PhysRevD.109.024064
% \href{http://arxiv.org/abs/arXiv:2212.14697}{arXiv:2212.14697}
% %5 citations counted in INSPIRE as of 25 Jun 2024
% \hfill  \underline{12 citations}

% %\cite{Tambalo:2022wlm}
% \item[\Num] \textit{``Gravitational wave lensing as a probe of halo properties and dark matter,''}
% G.~Tambalo, M.~Zumalac\'arregui, L.~Dai and M.~H.~Y.~Cheung, \\
% Phys. Rev. D \textbf{108}, no.10, 103529 (2023)
% % doi:10.1103/PhysRevD.108.103529
% \href{http://arxiv.org/abs/arXiv:2212.11960}{arXiv:2212.11960}
% %28 citations counted in INSPIRE as of 25 Jun 2024
% \hfill  \underline{49 citations}

% %\cite{Tambalo:2022plm}
% \item[\Num] \textit{``Lensing of gravitational waves: Efficient wave-optics methods and validation with symmetric lenses,''}
% G.~Tambalo, M.~Zumalac\'arregui, L.~Dai and M.~H.~Y.~Cheung, \\
% Phys. Rev. D \textbf{108}, no.4, 043527 (2023)
% % doi:10.1103/PhysRevD.108.043527
% \href{http://arxiv.org/abs/arXiv:2210.05658}{arXiv:2210.05658}
% %21 citations counted in INSPIRE as of 20 Jun 2024
% \hfill  \underline{40 citations}

% %\cite{Oancea:2022szu}
% \item[\Num] \textit{``Frequency- and polarization-dependent lensing of gravitational waves in strong gravitational fields,''}
% M.~A.~Oancea, R.~Stiskalek and M.~Zumalac\'arregui, \\
% Phys. Rev. D \textbf{109}, no.12, 124045 (2024)
% % doi:10.1103/PhysRevD.109.124045
% \href{http://arxiv.org/abs/arXiv:2209.06459}{arXiv:2209.06459}
% \hfill  \underline{21 citations}


% \item[\Num] \textit{``New Horizons for Fundamental Physics with LISA''},
% K.~G.~Arun \textit{et al.} [LISA],
% Living Rev. Rel. \textbf{25}, no.1, 4 (2022) %(\href{https://rdcu.be/cQKDM}{link})
% \href{http://arxiv.org/abs/arXiv:2205.01597}{arXiv:2205.01597}
% %shortened URL is now available for you to share full-text access to your paper by using the following link: https://rdcu.be/cQKDM
%  \quad (endorser) \hfill \underline{242 citations} \signhot{}\;\signcited{}

% %\cite{LISACosmologyWorkingGroup:2022jok}
% % \bibitem{LISACosmologyWorkingGroup:2022jok}
%  \item[\Num] \textit{``Cosmology with the Laser Interferometer Space Antenna''},
% P.~Auclair \textit{et al.} [LISA Cosmology Working Group],
% Living Rev. Rel. \textbf{26}, no.1, 5 (2023)
% \href{http://arxiv.org/abs/arXiv:2204.05434}{arXiv:2204.05434}
% \\ \qquad (section coordinator ``gravitational lensing'')  \hfill \underline{391 citations}  \signhot{}\;\signcited{}

% %\cite{Abdalla:2022yfr}
% % \bibitem{Abdalla:2022yfr}
%  \item[\Num] \textit{``Cosmology intertwined: A review of the particle physics, astrophysics, and cosmology associated with the cosmological tensions and anomalies''},
% E.~Abdalla \textit{et al.}
% JHEAp \textbf{34} (2022), 49-211
%  \href{http://arxiv.org/abs/arXiv:2203.06142}{arXiv:2203.06142} \quad (contributor)
% \hfill \underline{1,350 citations}  \signhot{}\;\signcited{}


% %\cite{Brando:2021jga}
% \item[\Num] \textit{``Fully relativistic predictions in Horndeski gravity from standard Newtonian N-body simulations,''}
% G.~Brando, K.~Koyama, D.~Wands, M.~Zumalac\'arregui, I.~Sawicki and E.~Bellini, \\
% JCAP \textbf{09}, 024 (2021)
% \href{http://arxiv.org/abs/arXiv:2105.04491}{arXiv:2105.04491}
% \hfill  \underline{10 citations}
% %6 citations counted in INSPIRE as of 25 Jun 2024

% %\cite{Traykova:2021hbr}
% \item[\Num] \textit{``Theoretical priors in scalar-tensor cosmologies: Shift-symmetric Horndeski models,''}
% D.~Traykova, E.~Bellini, P.~G.~Ferreira, C.~Garc\'\i{}a-Garc\'\i{}a, J.~Noller and M.~Zumalac\'arregui, \\
% Phys. Rev. D \textbf{104}, no.8, 083502 (2021)
% \href{http://arxiv.org/abs/arXiv:2103.11195}{arXiv:2103.11195}
% \hfill  \underline{38 citations}
% %33 citations counted in INSPIRE as of 26 Jun 2024

% %\cite{Ezquiaga:2020dao}
% \item[\Num]
% \textit{``Gravitational wave lensing beyond general relativity: birefringence, echoes and shadows''}
% \\{}J.~M.~Ezquiaga and M.~Zumalacarregui,
% \\ \href{http://arxiv.org/abs/arXiv:2009.12187}{arXiv:2009.12187}
% Phys. Rev. D \textbf{102}, no.12, 124048 (2020)
% \textbf{Editors' Suggestion}
% % doi:10.1103/PhysRevD.102.124048
% % [arXiv:2009.12187 [gr-qc]].
% %16 citations counted in INSPIRE as of 18 Jun 2021
% \hfill  \underline{77 citations} 


% %\cite{Zumalacarregui:2020cjh}
% \item[\Num]
% \textit{``Gravity in the Era of Equality: Towards solutions to the Hubble problem without fine-tuned initial conditions''} M.~Zumalacarregui,\\
% % doi:10.1103/PhysRevD.102.023523
% \href{http://arxiv.org/abs/arXiv:2003.06396}{arXiv:2003.06396}
% {Phys. Rev. D \textbf{102}, no.2, 023523 (2020)}
% % [arXiv:2003.06396 [astro-ph.CO]].
% %39 citations counted in INSPIRE as of 15 Jun 2021
% \hfill  \underline{98 citations}%$^\clubsuit$ 


% \item[\Num]
% \textit{``Prospects for Fundamental Physics with LISA,''}
% E.~Barausse %, E.~Berti, T.~Hertog, % S.~A.~Hughes, P.~Jetzer, P.~Pani, T.~P.~Sotiriou, N.~Tamanini, H.~Witek and K.~Yagi, 
% \textit{et al.}
% \href{http://arxiv.org/abs/arXiv:2001.09793}{arXiv:2001.09793}
% Gen. Rel. Grav. \textbf{52}, no.8, 81 (2020)
% % doi:10.1007/s10714-020-02691-1
% % [arXiv:2001.09793 [gr-qc]].
% \quad (endorser) \hfill  \underline{365 citations} 
% %76 citations counted in INSPIRE as of 18 Jun 2021



% %\cite{Garcia-Garcia:2019cvr}
% \item[\Num]
% \textit{``Theoretical priors in scalar-tensor cosmologies: Thawing quintessence''}
% \\{}C.~Garc\'\i{}a-Garc\'\i{}a, E.~Bellini, P.~G.~Ferreira, D.~Traykova and M.~Zumalacarregui,
% \\
% \href{http://arxiv.org/abs/arXiv:1911.02868}{arXiv:1911.02868}
% Phys. Rev. D \textbf{101}, no.6, 063508 (2020)
% \hfill  \underline{13 citations} 
% % doi:10.1103/PhysRevD.101.063508
% % [arXiv:1911.02868 [astro-ph.CO]].
% %5 citations counted in INSPIRE as of 16 Jun 2021


% %\cite{Bellini:2019syt}
% \item[\Num]
% \textit{``hi\_class: Background Evolution, Initial Conditions and Approximation Schemes''}
% \\{}E.~Bellini, I.~Sawicki and M.~Zumalacarregui,
% \\
% \href{http://arxiv.org/abs/arXiv:1909.01828}{arXiv:1909.01828}
% JCAP \textbf{02}, 008 (2020)
% \hfill  \underline{59 citations} 
% % doi:10.1088/1475-7516/2020/02/008
% % [arXiv:1909.01828 [astro-ph.CO]].
% %21 citations counted in INSPIRE as of 15 Jun 2021


%  \item[\Num] \textit{``Unveiling the Gravitational Universe at $\mu$-Hz Frequencies''},
% Sesana \textit{et al.},  %(w. M.~Zumalacarregui)
% Exper. Astron. \textbf{51}, no.3, 1333-1383 (2021)
%  \href{http://arxiv.org/abs/arXiv:1908.11391}{arXiv:1908.11391} 
%  %(forecasted constraints on graviton mass and gravitational lensing) \\ %, White Paper submitted to Voayage-2050
% \quad (contributor) \hfill  \underline{244 citations} \signcited{}

% \item[\Num] \textit{``High angular resolution gravitational wave astronomy''},
% Baker \textit{et al.},  %(w. M.~Zumalacarregui)
% Exper. Astron. \textbf{51}, no.3, 1441-1470 (2021)
%  \href{http://arxiv.org/abs/arXiv:1908.11410}{arXiv:1908.11410} %, White Paper submitted to the Voayage-2050
% \quad (contributor) \hfill  \underline{59 citations} 
 

% %\cite{Belgacem:2019pkk}
% % \bibitem{Belgacem:2019pkk}
% \item[\Num]
% \textit{``Testing modified gravity at cosmological distances with LISA standard sirens''}\\
%   E.~Belgacem {\it et al.} [LISA Cosmology Working Group] (1st tier author),\\
%   \href{http://arxiv.org/abs/arXiv:1906.0159}{arXiv:1906.0159} {JCAP {\bf 1907} (2019) 024} \hfill  \underline{210 citations} 

% \item[\Num]
% % \bibitem{Garcia-Garcia:2019ees}
% \textit{``$\alpha$-attractor dark energy in view of next-generation cosmological surveys''}\\
%   C.~Garc\'ia-Garc\'ia, P.~Ru\'iz-Lapuente, D.~Alonso and M.~Zumalacarregui,\\
%   \href{http://arxiv.org/abs/arXiv:1905.03753}{arXiv:1905.03753} {JCAP {\bf 1907} (2019) 025} \hfill  \underline{11 citation} 


% \item[\Num] %\cite{Brush:2018dhg}
% % \bibitem{Brush:2018dhg} 
% \textit{``No Slip CMB''},\\
%   M.~Brush, E.~V.~Linder and M.~Zumalacarregui, \\
%   \href{http://arxiv.org/abs/arXiv:1810.12337}{arXiv:1810.12337} 
%    {JCAP01 (2019) 029}
%   \hfill \underline{20 citations} 
% %   arXiv:1810.12337 [astro-ph.CO].
%   %%CITATION = ARXIV:1810.12337;%%

% \item[\Num] \textit{``Dark Energy in light of Multi-Messenger Gravitational-Wave astronomy''} \\
%  J. M. Ezquiaga, \textbf{M.~Zumalac\'arregui} \\
%   \href{https://arxiv.org/abs/1807.09241}{arXiv:1807.09241} %\\
%  \textbf{Invited review},  {Front. Astron. Space Sci. (2018) 44} \hfill \underline{205 citations} 
%  %: ``Dark Energy and the Acceleration of the Universe''
 
%  \item[\Num] \textit{``Dark energy from $\alpha$-attractors: phenomenology and observational constraints''}
% C. Garc\'ia-Garc\'ia, E. Linder, P. Ru\'iz-Lapuente, \textbf{M.~Zumalac\'arregui}\\
%  \href{https://arxiv.org/abs/1803.00661}{arXiv:1803.00661} 
%  {JCAP08 (2018) 022}
% \hfill \underline{48 citations} 


%  \item[\Num] \textit{``Testing (modified) gravity with 3D and tomographic cosmic shear''},  \\
%  A. Spurio Mancini , R. Reischke, V. Pettorino, B.M. Sch\"afer and \textbf{M.~Zumalac\'arregui}  \\
%   \href{https://arxiv.org/abs/1801.04251}{arXiv:1801.04251} 
%   {Mon.Not.Roy.Astron.Soc. 480 (2018) 3725}
%  \hfill  \underline{55 citations} 

%  \item[\Num] \textit{``Limits on stellar-mass compact objects as dark matter from gravitational lensing of type Ia supernovae''}, 
%  \textbf{M.~Zumalac\'arregui} and U.~Seljak, \\
%   \href{https://arxiv.org/abs/1712.02240}{arXiv:1712.02240} 
%   {Phys. Rev. Lett. 121 (2018) 141101} 
%    \\ (Editors' Suggestion and \href{https://physics.aps.org/articles/v11/97}{APS Physics Viewpoint})
%  \hfill  \underline{269 citations} \signcited 
 
%  %\cite{Ezquiaga:2017ekz}
% % \bibitem{Ezquiaga:2017ekz}
%  \item[\Num] \textit{``Dark Energy after GW170817: dead ends and the road ahead''} \\
%  J.~M.~Ezquiaga and \textbf{M.~Zumalac\'arregui}, \\
%  \href{https://arxiv.org/abs/1710.05901}{arXiv:1710.05901} 
%  {Phys.Rev.Lett. 119 (2017) 251304}
%  \\ (Editors' Suggestion and \href{https://physics.aps.org/articles/v10/134}{APS Physics Viewpoint})
%  \hfill  \underline{936 citations} \signcited
% %   arXiv:1710.05901 [astro-ph.CO].
%   %%CITATION = ARXIV:1710.05901;%%
%   %1 citations counted in INSPIRE as of 18 Oct 2017
  
%   %\cite{Bellini:2017avd}
% \item[\Num] \textit{``Comparison of Einstein-Boltzmann solvers for testing General Relativity''}
%   \\ E. Bellini, A. Barreira, N. Frusciante, B. Hu, S. Peirone, M. Raveri, \textbf{M. Zumalac\'arregui}, {\it et al.},\\
%   \href{https://arxiv.org/abs/1709.09135}{arXiv:1709.09135} \textbf{Phys.Rev. D97 (2018) 023520}
%   \hfill  \underline{77 citations}
%   %%CITATION = ARXIV:1709.09135;%%
%   %1 citations counted in INSPIRE as of 29 Sep 2017
  
% \item[\Num]  \textit{``Galileon Gravity in Light of ISW, CMB, BAO and $H_0$ data''} \\
% J.~Renk, M.~Zumalacarregui, F.~Montanari and A.~Barreira, \\
%  \href{http://arxiv.org/abs/arXiv:1707.02263}{arXiv:1707.02263} \textbf{JCAP10 (2017) 020}
%  \hfill  \underline{193 citations}
% %(Jul 7, 2017)
% % \inspireurl{http://inspirehep.net/record/1609059}

% % \bibitem{Ezquiaga:2017ner}
% \item[\Num] \textit{``Field redefinitions in theories beyond Einstein gravity using the language of differential forms''} \\
%   J.~M.~Ezquiaga, J.~Garc\'ia-Bellido and \textbf{M.~Zumalac\'arregui},\\
%   \href{http://arxiv.org/abs/arXiv:1701.05476}{arXiv:1701.05476} \textbf{Phys.Rev. D95 (2017) 084039} 
%   \hfill  \underline{29 citations}
  

% \item[\Num]  \textit{``Observational Future of Cosmological Scalar-Tensor Theories''}\\
%   D.~Alonso, E.~Bellini, P.~G.~Ferreira, \textbf{M.~Zumalac\'arregui} \\
%   \href{http://arxiv.org/abs/arXiv:1610.09290}{arXiv:1610.09290} {Phys.Rev. D95 (2017) no.6, 063502} 
%   \hfill \underline{151 citations} % 
  
%     %\cite{Bettoni:2016mij} 
% \item[\Num]  \textit{``Speed of Gravitational Waves and the Fate of Scalar-Tensor Gravity''}\\
%   D.~Bettoni, J.~M.~Ezquiaga, K.~Hinterbichler, \textbf{M.~Zumalac\'arregui} \\
%   \href{http://arxiv.org/abs/arXiv:1608.01982}{arXiv:1608.01982}  {Phys.Rev. D95 (2017) no.8, 084029} 
%   \hfill  \underline{192 citations}%$^\clubsuit$

% %\cite{Konnig:2016idp}
% \item[\Num] 
% \textit{``A spectre is haunting the cosmos: Quantum stability of massive gravity with ghosts''}\\
%  F.~K\"onnig, H.~Nersisyan, Y.~Akrami, L.~Amendola, \textbf{M.~Zumalac\'arregui} \\
%   \href{http://arxiv.org/abs/arXiv:1605.08757}{arXiv:1605.08757} {JHEP11 (2016) 118} \hfill 
%   \hfill \underline{15 citations}
    
%     %\cite{Zumalacarregui:2016pph}
% \item[\Num]  \textit{``hi\_class: Horndeski in the Cosmic Linear Anisotropy Solving System''}\\
%  \textbf{M.~Zumalac\'arregui}, E.~Bellini, I.~Sawicki, J.~Lesgourgues and P.~Ferreira \\
%   \href{http://arxiv.org/abs/arXiv:1605.06102}{arXiv:1605.06102} {JCAP08 (2017) 019} 
%   \hfill \underline{205 citations}   

  
% \item[\Num] \textit{``Gravity at the horizon: on relativistic effects, CMB-LSS correlations and ultra-large scales in Horndeski theory''}
% J.~Renk, \textbf{M.~Zumalac\'arregui}, F.~Montanari \\
%   \href{http://arxiv.org/abs/1603.01269}{arXiv:1604.03487} {JCAP07 (2016) 040} 
%   \hfill  \underline{59 citations}
% % \href{http://inspirehep.net/record/XXXX}{HEP entry}

% \item[\Num] \textit{``Towards the most general scalar-tensor theories of gravity: a unified approach in the language of differential forms''} %\\
%   J.~M.~Ezquiaga, J.~Garc\'ia-Bellido and \textbf{M.~Zumalac\'arregui} \\
%   \href{http://arxiv.org/abs/1603.01269}{arXiv:1603.01269} {Phys. Rev. D 94 (2016) 024005}. 
%   \hfill \underline{47 citations}

% \item[\Num] 
% \textit{ ``Non-linear evolution of the BAO scale in alternative theories of gravity''} \\
% E. Bellini,  M.~Zumalacarregui, \\
%   \href{http://arxiv.org/abs/1505.03839}{arXiv:1505.03839}. {Phys.Rev. D92 (2015)}.
%  \hfill \underline{38 citations} % [astro-ph.CO]   
  
% \item[\Num] 
% \textit{``Surfing gravitational waves: can bigravity survive growing tensor modes?''} \\
% L. Amendola, F. K\"onnig, M. Martinelli, V. Pettorino, M.~Zumalacarregui, \\
%   \href{http://arxiv.org/abs/1503.02490}{arXiv:1503.02490}. {JCAP05 (2015) 052}.
%  \hfill \underline{36 citations} % [astro-ph.CO] 

% \item[\Num]
% \textit{``Shaken, not stirred: kinetic mixing in scalar-tensor theories of gravity''} \\
% D. Bettoni, M.~Zumalacarregui, \\
%   \href{http://arxiv.org/abs/1502.02666}{arXiv:1502.02666} {Phys.Rev. D91 (2015) 104009}. 
%  \hfill  \underline{56 citations} % [astro-ph.CO] 

%     \item[\Num]
%  \textit{``Transforming gravity: from derivative couplings to matter 
%  to second-order scalar-tensor theories beyond the Horndeski Lagrangian''} \\ % $^\star$
%  M.~Zumalacarregui, J. Garcia-Bellido, \\
% \href{http://arxiv.org/abs/arXiv:1308.4685}{arXiv:1308.4685}. {Phys.Rev. D89 (2014) 064046}. 
% \hfill \underline{559 citations} \signcited % [astro-ph.CO] 

%  \item[\Num]
%  \textit{``DBI Galileons in the Einstein Frame: Local Gravity and Cosmology''} \\ % $^\star$
%  M.~Zumalacarregui, T. Koivisto and D. Mota, \\
% \href{http://arxiv.org/abs/arXiv:1210.8016}{arXiv:1210.8016}, {Phys.Rev. D87 (2013) 083010}.  
% \hfill \underline{149 citations} % [astro-ph.CO] 

% \item[\Num] 
% \textit{``Screening Modifications of Gravity through Disformally Coupled Fields''} \\ % $^\star$
%  T. Koivisto, D. Mota, and M.~Zumalacarregui, \\
%  \href{http://arxiv.org/abs/arXiv:1205.3167}{arXiv:1205.3167}, {Phys.Rev.Lett. 109 (2012) 241102}. 
% \hfill \underline{146 citations} % [astro-ph.CO] 

% \item[\Num] 
% \textit{``Tension in the Void: Cosmic Rulers Strain Inhomogeneous Cosmologies''} \\ % $^\star$
%  M.~Zumalacarregui, J. Garcia-Bellido and P. Ruiz-Lapuente. \\
%  \href{http://arxiv.org/abs/arXiv:1201.2790}{arXiv:1201.2790}. {JCAP10 (2012) 009}. 
% \hfill \underline{57 citations} % [astro-ph.CO] 

% \item[\Num] \textit{``Constraining Entropic Cosmology''} \\ %think about it. $^\star$
%  T. Koivisto, D. F. Mota and M.~Zumalacarregui \\
%  \href{http://arxiv.org/abs/arXiv:1011.2226}{arXiv:1011.2226}, {JCAP02 (2011) 027}. 
% \hfill \underline{36 citations} % [astro-ph.CO] 

%  \item[\Num] \textit{``Disformal Scalar Fields and the Dark Sector of the Universe''} \\ %$^\star$
%   M.~Zumalacarregui, T. Koivisto, D. F. Mota and P. Ruiz-Lapuente, \\
%  \href{http://arxiv.org/abs/arXiv:1004.2684}{arXiv:1004.2684}, {JCAP05 (2010) 038}. 
% \hfill \underline{111 citations} % [astro-ph.CO] 
 
% \end{itemize}


% % \subsection*{White Papers}
% % \begin{itemize}\setlength{\itemsep}{0.5pt}


% % \end{itemize}


% % \subsection*{Conference proceedings}%Scientific Publications %\LARGE 
% % %\href{http://inspirehep.net/search?p=author\%3AM.Zumalacarregui.1}
% % \begin{itemize}\setlength{\itemsep}{0.pt}
% % \item[\Num] \textit{``No LIGO MACHO: size matters, the conclusion stands''}, 
% % M.~Zumalacarregui \\ Rencontres de Moriond 2018. %\\
% % %  \href{http://arxiv.org/abs/arXiv:1202.1281}{arXiv:1202.1281}, AIP Conf.Proc. 1458 (2011) 539-542.
% % \item[\Num] \textit{``Modified Entropic Gravity and Cosmology''}, 
% % M.~Zumalacarregui, \\
% %  \href{http://arxiv.org/abs/arXiv:1202.1281}{arXiv:1202.1281}, AIP Conf.Proc. 1458 (2011) 539-542.
% % \end{itemize}

% \subsection*{Preprints}\setlength{\itemsep}{0.pt}
% \begin{itemize}
% \item[\Num] \textit{``Black holes as telescopes: Discovering supermassive binaries through quasi-periodic lensed starlight''} H.~Wang, M.~Zumalacarregui, B.~Kocsis
% \href{https://arxiv.org/abs/2506.16544}{2506.16544}

% \item[\Num] \textit{``To the Problem of Cosmic Expansion in Massive Gravity,''}  
% L.~Heisenberg, A.~Longo, G.~Tambalo and M.~Zumalacarregui,  
% \href{http://arxiv.org/abs/arXiv:2411.19873}{arXiv:2411.19873}.  



% \item[\Num] \textit{``wolensing: A Python package for computing the amplification factor for gravitational waves with wave-optics effects,''}  
% S.~M.~C.~Yeung, M.~H.~Y.~Cheung, M.~Zumalacarregui and O.~A.~Hannuksela,  
% \href{http://arxiv.org/abs/arXiv:2410.19804}{arXiv:2410.19804}.  
%  \hfill  \underline{1 citation} 


% %\cite{Euclid:2024yrr}
% \item[\Num] \textit{``Euclid. I. Overview of the Euclid mission,''}
% Y.~Mellier \textit{et al.} [Euclid],
% \href{http://arxiv.org/abs/arXiv:2405.13491}{arXiv:2405.13491}.
% %14 citations counted in INSPIRE as of 26 Jun 2024
%  \hfill  \underline{166 citations} 

% %\cite{Zumalacarregui:2024ocb}
% \item[\Num] \textit{``Lens Stochastic Diffraction: A Signature of Compact Structures in Gravitational-Wave Data,''}
% M.~Zumalacarregui,
% \href{http://arxiv.org/abs/arXiv:2404.17405}{arXiv:2404.17405}.
% [arXiv:2404.17405 [gr-qc]].
%  \hfill  \underline{6 citations} 

% %1 citations counted in INSPIRE as of 25 Jun 2024

% % \bibitem{Kalogera:2021bya}
%  \item[\Num] \textit{`The Next Generation Global Gravitational Wave Observatory: The Science Book''},
% V.~Kalogera \textit{et al.} %B.~S.~Sathyaprakash, M.~Bailes, M.~A.~Bizouard, A.~Buonanno, A.~Burrows, M.~Colpi, M.~Evans, S.~Fairhurst and S.~Hild,
%  \href{http://arxiv.org/abs/arXiv:2111.06990}{arXiv:2111.06990}
% \qquad (contributed Fig.~5) 
% \hfill \underline{121 citations} 

% \item[\Num] \textit{``Extreme Gravity and Fundamental Physics''},
% Sathyaprakash \textit{et al.}, % (w. M.~Zumalacarregui)
%  \href{http://arxiv.org/abs/arXiv:1903.09221}{arXiv:1903.09221}, 
%  % White Paper submitted to the Astro-2020 
% (endorser)\\ \phantom{z} \hfill  \underline{71 citations} 
 

% \end{itemize}


\section*{Professional References}
\begin{itemize}
\item Prof. \textbf{Alessandra Buonanno} (MPI-Grav Potsdam, ACR Division Director) - \url{buonanno@aei.mpg.de}
\item Prof. \textbf{Matias Zaldarriaga} (IAS) - \url{matiasz@ias.edu}
 \item Prof. \textbf{Pedro G. Ferreira} (Oxford, Head of Astrophysics) - \url{pedro.ferreira@physics.ox.ac.uk}
\item Prof. \textbf{Uro\v{s} Seljak} (UC Berkeley \& LBL, BCCP Director) - \url{useljak@berkeley.edu}
% \item Prof. \textbf{Diego Blas} (ICREA-IFAE) - \url{dblas@ifae.es}
\item Dr. \textbf{Filippo Vernizzi} (CEA, IPhT Saclay) - \url{filippo.vernizzi@ipht.fr }
 \item Prof. \textbf{Julien Lesgourgues} (RWTH Aachen University) - \url{lesgourg@physik.rwth-aachen.de}
 \item Prof. \textbf{Holger M\"uller} (UC Berkeley) - \url{hm@berkeley.edu}
 \item Prof. \textbf{Luca Amendola} (University of Heidelberg) - \url{amendola@thphys.uni-heidelberg.de}
% \item Prof. \textbf{Pilar Ruiz-Lapuente} (IFF-CSIC) - \url{pilarrl@iff.csic.es}
%  \item Prof. \textbf{David F. Mota} (Oslo) - \url{d.f.mota@astro.uio.no}
%  \item Dr. \textbf{Tomi S. Koivisto} (Nordita) - \url{tomi.koivisto@nordita.org}
 \item Prof. \textbf{Juan Garcia-Bellido} (IFT-UAM/CSIC) - \url{juan.garciabellido@uam.es}
\end{itemize}

\end{document}

% \newpage


% \newpage

% % \newpage
% \noindent
% Collaborations
% \vspace{-0.3cm}
% \begin{longtable}{c p{15cm}}
%  2017- &  Collaboration with \emph{Sparrho} popular sicence website %\\ &
%    \href{https://www.sparrho.com/u/2b878/}{\tt www.sparrho.com}
%    \\[4pt]
%  2017- &  Collaboration with the \emph{FdeT} popular sicence website (in Spanish) \\ & \url{http://fdetonline.com/author/miguel-zumalacarregui}
%   \\[4pt]
%  2016 & \textbf{Casting consultant} in the movie \textit{``Las leyes de la Termodinamica''} by Mateo Gil \\[4pt]
%  2015- &   \textbf{Blog:} {\it ``Persistencia y Serendipia''}
%   (in Spanish) \, \url{www.fronterad.com/?q=blog/2050} 
% \end{longtable}
% % \vspace{-0.3cm}


\newpage
